\section{임피던스 제어}

로봇 손끝이 목표 위치에 정확히 도착하는 것만으로는 접촉 작업을 설명하기 어렵다.
벽, 부품, 사람의 손을 만나는 순간에는 위치 오차가 곧 힘이 되고, 너무 단단한 제어는 작은 오차도 큰 충격으로 바꿀 수 있다.
이 장은 그 한계에서 출발해, 질량-스프링-댐퍼 모델을 로봇 끝단의 목표 동역학으로 읽는 방법을 다룬다.
평형점, 전달함수, 관절/직교 임피던스, 가상 벽을 차례로 보면서 강성, 감쇠, 관성이 플롯에서 어떤 흔적으로 나타나는지 확인한다.

이 장은 세 장면으로 읽으면 쉽다.
자유공간에서는 목표점으로 얼마나 빠르고 안정적으로 가는지를, 접촉 뒤 완전히 멈춘 장면에서는 힘과 위치 오차의 정적 관계를, 닿는 순간부터 멈출 때까지는 진동, 오버슈트, 감쇠, 관성을 본다.
같은 식이라도 어느 장면을 보느냐에 따라 주인공이 달라진다.

\begin{table}[!htbp]
  \centering
  \caption{임피던스 제어를 읽는 세 장면}
  \small
  \setlength{\tabcolsep}{4pt}
  \begin{tabularx}{\linewidth}{@{}>{\raggedright\arraybackslash}p{0.24\linewidth}>{\raggedright\arraybackslash}p{0.31\linewidth}>{\raggedright\arraybackslash}X@{}}
    \toprule
    장면 & 먼저 보는 질문 & 주로 드러나는 파라미터 \\
    \midrule
    자유공간 이동 & 목표점까지 얼마나 빠르고 매끄럽게 가는가? & 고유진동수, 감쇠비, 포화, 목표 궤적 \\
    정적 접촉 & 오래 기다린 뒤 얼마나 밀려 있고 힘은 얼마인가? & 강성, 평형점, 목표 오프셋 \\
    과도 접촉 & 닿는 순간 얼마나 튀고 흔들리며 멈추는가? & 관성, 감쇠, 속도 제한, 환경 강성 \\
    \bottomrule
  \end{tabularx}
\end{table}

\subsection{왜 임피던스 제어가 등장했는가?}

전통적인 로봇 제어를 단순화하면, 로봇에게 어디로 가야 하는지를 알려주고 그 위치를 정확히 따라가게 만드는 방식이었다.
공장 안에서 같은 부품을 같은 위치에 반복해서 옮기는 작업에서는 이 방식만으로도 충분히 강력하다.
대상이 단단히 고정되어 있고, 작업 경로가 미리 정해져 있으며, 사람과 환경의 예기치 못한 접촉이 적다면 위치 정확도가 가장 중요한 목표가 된다.

하지만 조작(manipulation)은 단순 이동보다 어렵다.
로봇 손끝이 문고리를 잡거나, 표면을 닦거나, 부품을 구멍에 끼우거나, 사람이 건네는 물체를 받는 순간에는 위치만 맞추는 것으로 충분하지 않다.
작은 위치 오차가 큰 힘으로 바뀔 수 있고, 그 힘은 물체를 망가뜨리거나 로봇을 불안정하게 만들 수 있다.
반대로 힘을 너무 약하게 만들면 작업을 수행하지 못한다.
이때 로봇은 위치와 힘을 따로따로 다룰 수 없고, 둘 사이의 관계를 설계해야 한다.
예를 들어 구멍에 부품을 끼우는 상황을 생각해보자.
옆 방향으로는 위치를 잘 맞추어야 하지만, 삽입 방향으로는 너무 단단하게 밀어붙이면 부품이 걸리거나 망가질 수 있다.
반대로 삽입 방향을 너무 느슨하게 두면 부품이 들어가지 않는다.
이런 작업에서는 모든 방향에서 위치를 세게 고정하는 것도, 모든 방향에서 힘만 맞추는 것도 답이 아니다.
방향마다 어느 정도 단단하고 어느 정도 부드럽게 반응할지를 정해야 한다.

네빌 호건(Neville Hogan)이 1985년에 제시한 임피던스 제어의 핵심도 여기에 있다 \cite{hogan1985impedancepart1}.
로봇이 환경과 접촉할 때 중요한 것은 특정 위치나 특정 힘 하나를 고집하는 것이 아니라, 외부 힘과 로봇 운동 사이의 동역학적 관계이다.
호건의 관점 전환은 위치냐 힘이냐를 먼저 고르는 데 있지 않았다.
접촉점에서 힘과 운동이 어떤 관계로 오가야 하는지를 먼저 설계하자는 데 핵심이 있다.
접촉 제어를 설명할 때 흔히 함께 떠올리는 배경에는 방향별로 위치를 제어할지 힘을 제어할지 나누어 생각하는 하이브리드 위치/힘 제어가 있다.
임피던스 제어는 그 선택 이전에 접촉 포트의 힘-운동 관계를 설계한다.
사람의 팔을 생각해보면 직관적이다.
손을 책상 위에 올릴 때 우리는 손 위치를 완전히 고정하지도 않고, 힘만 무작정 일정하게 만들지도 않는다.
팔은 어느 정도 부드럽게 밀리고, 빠른 충격에는 감쇠를 보이며, 필요할 때는 단단하게 버틴다.
임피던스 제어는 로봇에도 이런 상호작용의 성질을 설계하려는 접근이다.

\subsection{기본 아이디어}

임피던스 제어(impedance control)는 로봇에게 목표 위치를 무조건 정확히 따라가라고 명령하는 방식이 아니다.
대신 로봇 끝단이 외부 힘을 받았을 때, 마치 우리가 설계한 가상의 질량-스프링-댐퍼처럼 반응하도록 만드는 제어 방법이다 \cite{hogan1985impedancepart1}.
질량-스프링-댐퍼 장의 방정식은 자연에 존재하는 물체의 운동을 설명하는 모델이었다.
여기서는 같은 2차 시스템을 로봇이 일부러 따르게 만들 목표 동역학으로 사용한다.
로봇 안에 진짜 스프링이나 댐퍼를 넣는다는 뜻이 아니라, 제어기가 그런 물리적 감각의 힘과 운동 관계를 만들겠다는 약속이다.
이 식이 답하려는 질문은 세 가지이다.
외부 힘을 받았을 때 얼마나 밀릴 것인가, 얼마나 빠르게 돌아올 것인가, 돌아오는 동안 얼마나 흔들릴 것인가.

다만 목표 동역학을 종이에 쓴다고 실제 로봇이 자동으로 그대로 움직이는 것은 아니다.
아래 식은 우리가 로봇에 기대하는 이상적인 힘-운동 관계이고, 실제 구현에서는 그 관계에 가까워지도록 모터 토크나 목표 위치를 계속 계산한다.
이 해석이 잘 맞으려면 몇 가지 전제가 필요하다.
\begin{itemize}
  \item 로봇이 필요한 힘이나 토크를 낼 수 있어야 한다.
  \item 중력, 관성, 마찰 같은 로봇 자체 동역학을 어느 정도 보상할 수 있어야 한다.
  \item 센서와 제어 주기가 충분히 빨라서 접촉 변화에 늦게 반응하지 않아야 한다.
  \item 모터 포화, 지연, 큰 모델 오차가 목표 임피던스를 압도하지 않아야 한다.
\end{itemize}
제어 주기는 제어기가 센서를 읽고 명령을 다시 보내는 시간 간격이고, 포화는 모터가 낼 수 있는 힘이나 토크의 한계에 걸린 상태이다.
대역폭은 로봇이 얼마나 빠른 변화까지 따라갈 수 있는지를 나타내는 말로, 느린 제어기는 빠른 접촉 변화를 뒤늦게 따라가게 된다.
즉 임피던스 제어식은 현실을 생략한 공식이 아니라, 충분히 빠르고 강한 제어기가 실제 로봇의 응답을 가상의 물리 시스템에 가깝게 보이도록 만드는 설계 목표이다.

가장 단순한 1차원 임피던스 모델은 다음과 같다.
오른쪽의 $f_{\mathrm{ext}}$는 환경이 로봇에 가하는 부호 있는 힘이다.
\begin{equation}
  m_d(\ddot{x}-\ddot{x}_d)
  +
  d_d(\dot{x}-\dot{x}_d)
  +
  k_d(x-x_d)
  =
  f_{\mathrm{ext}}.
  \label{eq:impedance_1d}
\end{equation}
여기서 $x$는 실제 위치, $x_d$는 희망 위치, $m_d$는 희망 관성, $d_d$는 원하는 감쇠계수, $k_d$는 희망 강성이다.
$m_d$는 로봇 링크의 실제 질량이 아니라, 로봇 끝단이 그렇게 느껴지도록 만들고 싶은 가상 관성이다.
실제 링크 질량, 자세에 따라 달라지는 끝단 유효질량, 제어기가 설계하는 희망 관성은 뒤에서 따로 구분한다.
$f_{\mathrm{ext}}$는 외부에서 로봇에 가해지는 힘이다.
로봇이 환경에 닿거나 사람이 로봇을 밀면 이 외력이 생긴다.

앞으로 힘의 방향을 헷갈리지 않도록 기준을 먼저 정리하자.
$f_{\mathrm{ext}}$는 환경이 로봇에 가하는 힘을 방향까지 포함해 쓴 값이다.
$f_{\mathrm{cmd}}$는 로봇 제어기가 만들려고 하는 명령힘이다.
$f_{\mathrm{wall}}$은 벽이 로봇에 가하는 힘이다.
반면 $f_{\mathrm{contact}}$라고 쓸 때는 문맥에 따라 접촉력의 크기를 말하는 경우가 많으므로, 크기를 강조할 때는 $|f_{\mathrm{contact}}|$처럼 쓴다.
즉, 로봇이 벽을 오른쪽으로 미는 힘과 벽이 로봇을 왼쪽으로 미는 힘은 크기는 같아도 부호는 반대이다.

3차원 작업공간의 병진 위치에 대해서는 같은 생각을 벡터와 행렬로 쓴다.
여기서 새로 어려운 물리를 추가하는 것은 아니다.
1차원에서 보던 위치 오차, 속도, 강성, 감쇠를 $x$, $y$, $z$ 세 방향으로 나누어 쓴 것부터 시작한다.
행렬은 방향별 스프링과 댐퍼를 한꺼번에 적기 위한 표기이다.
처음에는 1차원 식을 $x$, $y$, $z$ 세 방향에 나란히 붙인 것으로 생각하면 된다.
비대각 강성이나 행렬의 에너지 해석은 뒤의 가상 퍼텐셜 에너지 절에서 다시 볼 것이므로, 여기서는 축마다 스프링과 댐퍼를 하나씩 둔 그림부터 붙잡자.
회전, 자세, 렌치는 더 정확한 로봇 제어를 위해 필요하지만, 이 글에서는 먼저 손끝 위치가 어떻게 밀리고 되돌아오는지에 집중한다.
\begin{equation}
  \bm{M}_d(\ddot{\bm{x}} - \ddot{\bm{x}}_d)
  +
  \bm{D}_d(\dot{\bm{x}} - \dot{\bm{x}}_d)
  +
  \bm{K}_d(\bm{x} - \bm{x}_d)
  =
  \bm{f}_{\mathrm{ext}}.
  \label{eq:task_impedance}
\end{equation}
여기서 $\bm{x}\in\R^3$는 로봇 끝단의 병진 위치, $\bm{x}_d$는 희망 위치이다.
$\bm{M}_d$, $\bm{D}_d$, $\bm{K}_d$는 각각 희망 관성, 감쇠, 강성 행렬이다.
첫 항은 가속도 변화를 쉽게 허용하지 않는 가상 무거움이고, 둘째 항은 속도에 비례한 브레이크이며, 셋째 항은 위치 오차에 비례해 외력과 균형을 이루는 가상 스프링 항이다.
식 \eqref{eq:task_impedance}의 좌변에 있는 $\bm{K}_d(\bm{x}-\bm{x}_d)$는 그 오차를 유지하려면 외부가 넣어야 하는 힘 항으로 읽고, 제어기가 목표점으로 되돌리기 위해 계산하는 복원 명령은 뒤에서 $-\bm{K}_d(\bm{x}-\bm{x}_d)$처럼 반대 부호로 쓴다.
예를 들어 $\bm{K}_d=\diag(k_x,k_y,k_z)$라면 $x$, $y$, $z$ 방향에 서로 다른 가상 스프링을 단 것과 같다.
$k_x$를 작게 하고 $k_y$, $k_z$를 크게 하면 한 방향으로는 잘 밀리고 다른 방향으로는 단단하게 버티는 끝단을 만들 수 있다.
대각이 아닌 항이 있으면 한 방향의 오차가 다른 방향의 힘에도 섞인다는 뜻이다.
예를 들어 $K_{xy}\neq0$이면 $y$ 방향 오차가 $x$ 방향 복원력에도 영향을 주므로, 가상 스프링의 주된 방향이 좌표축과 비스듬히 놓인 것처럼 해석할 수 있다.
다만 에너지 지형으로 해석하려면 강성 행렬이 대칭이어야 하므로, 비대각 항을 쓸 때는 보통 $K_{xy}=K_{yx}$인 보존적인 가상 스프링부터 생각한다.
로봇에서 위치(position)는 보통 $x,y,z$ 좌표를 뜻한다.
자세(orientation)는 손끝이 어느 방향으로 돌아가 있는지를 뜻한다.
포즈(pose)는 위치와 자세를 합친 말이다.
힘(force)은 병진 방향으로 미는 작용이고, 토크(torque)는 회전시키는 작용이다.
렌치(wrench)는 힘과 토크를 한 벡터로 묶은 것이다.
따라서 자세까지 포함한 6차원 임피던스는 별도의 회전 오차 표현과 힘-토크를 합친 렌치를 사용해야 하므로, 이 글에서는 먼저 병진 위치 임피던스에 집중한다.
\vmark{scope-pose-wrench}%
이 단순화는 자세가 별도 제어기로 충분히 유지되거나, 접촉 모멘트를 잠시 무시해도 되는 교육용 장면에 맞춘 것이다.
문손잡이를 돌리거나 공구의 각도를 정밀하게 맞추는 작업처럼 자세와 모멘트가 작업 성공을 좌우하는 경우에는 6차원 렌치와 자세 임피던스를 함께 다루어야 한다.

교육용 데모에서는 먼저 병진 위치만 다루는 것이 좋다.
즉,
\begin{equation}
  \bm{x} =
  \begin{bmatrix}
  x & y & z
  \end{bmatrix}^{\mathsf{T}}
  \in \R^3
\end{equation}
로 두고, 자세 제어는 나중에 추가한다.
또한 목표점이 정지해 있다고 보면
\begin{equation}
  \dot{\bm{x}}_d = \bm{0}, \qquad \ddot{\bm{x}}_d = \bm{0}
\end{equation}
으로 단순화할 수 있다.
그러면 목표점 주변에서 발생하는 가상 스프링-댐퍼 힘은
\begin{equation}
  \bm{f}_{\mathrm{cmd}}
  =
  -\bm{K}_d(\bm{x}-\bm{x}_d)
  -
  \bm{D}_d\dot{\bm{x}}
  \label{eq:virtual_spring_damper}
\end{equation}
가 된다.
식 \eqref{eq:task_impedance}은 제어기까지 포함한 로봇이 외력에 대해 보여 주길 원하는 닫힌루프 행동이다.
반면 식 \eqref{eq:virtual_spring_damper}의 $\bm{f}_{\mathrm{cmd}}$는 그 행동에 가까워지도록 제어기가 위치 오차와 속도에서 계산하는 복원 힘이다.
두 식은 같은 종류의 식이 아니다.

\begin{table}[!htbp]
  \centering
  \caption{목표 관계와 구현 명령의 차이}
  \small
  \setlength{\tabcolsep}{4pt}
  \begin{tabularx}{\linewidth}{@{}>{\raggedright\arraybackslash}p{0.25\linewidth}>{\raggedright\arraybackslash}p{0.25\linewidth}>{\raggedright\arraybackslash}X@{}}
    \toprule
    구분 & 읽는 방향 & 의미 \\
    \midrule
    목표 임피던스 관계 & 외력 $\rightarrow$ 움직임 & 밖에서 힘을 받았을 때 로봇 전체가 얼마나 밀리고 흔들릴지 정하는 닫힌루프 목표이다. \\
    가상 스프링-댐퍼 명령 & 오차와 속도 $\rightarrow$ 복원힘 & 목표 관계에 가까워지도록 제어기가 매 순간 계산하는 내부 명령이다. \\
    \bottomrule
  \end{tabularx}
\end{table}

식 \eqref{eq:task_impedance}의 $\bm{f}_{\mathrm{ext}}$는 환경이 로봇에 가한 외력이다.
바깥에서 힘을 받았을 때 로봇과 제어기 전체가 얼마나 밀리고, 얼마나 빠르게 움직이고, 얼마나 가속하는지를 정하는 목표 관계이다.
반면 식 \eqref{eq:virtual_spring_damper}의 $\bm{f}_{\mathrm{cmd}}$는 제어기가 매 순간 계산해 목표점 쪽으로 되돌리려는 내부 명령힘이다.
같은 위치 오차 $e=x-x_d$라도, $f_{\mathrm{ext}}=k_de$는 로봇을 그 위치에 밀어 둔 외력으로 읽고, $f_{\mathrm{cmd}}=-k_de$는 그 외력에 맞서 목표점으로 돌아오려는 복원 명령으로 읽으면 부호 기준이 고정된다.
정지 평형에서는 속도와 가속도가 거의 0이므로, 이 두 힘은 같은 직선 위에서 반대 부호로 균형을 이룬다.
스칼라 예로는 환경이 $+F_0$로 로봇을 밀어 $e_{\mathrm{ss}}=F_0/k_d$만큼 밀어 두고, 제어기가 $-F_0$에 가까운 복원 명령을 만들어 버티는 그림이다.

이 관계가 실제 plant 위에서 어떻게 닫힌루프 식이 되는지도 1축으로 보면 어렵지 않다.
손끝의 한 방향을 등가 질량 $m_{\mathrm{eff}}$로 단순화하고, 제어기가 계산한 명령힘과 외력이 함께 작용한다고 하면
\begin{equation}
  m_{\mathrm{eff}}\ddot{x}=f_{\mathrm{cmd}}+f_{\mathrm{ext}}
\end{equation}
이다.
목표점이 정지해 있고 $e=x-x_d$라면 $\dot{e}=\dot{x}$, $\ddot{e}=\ddot{x}$이므로, 가상 스프링-댐퍼 명령
\begin{equation}
  f_{\mathrm{cmd}}=-d_d\dot{e}-k_de
\end{equation}
를 대입해
\begin{equation}
  m_{\mathrm{eff}}\ddot{e}+d_d\dot{e}+k_de=f_{\mathrm{ext}}
\end{equation}
를 얻는다.
즉 실제 plant의 관성이 그대로 남고, 제어기가 만든 스프링항과 댐퍼항이 닫힌루프의 위치항과 속도항을 만든다.
만약 원하는 식의 질량 $m_d$까지 실제 $m_{\mathrm{eff}}$와 다르게 만들고 싶다면, 단순한 스프링-댐퍼 명령만으로는 부족하고 동역학 보상이나 관성 shaping이 필요하다.
위치 명령형 인터페이스에서는 이 관성 shaping을 엄밀히 구현하기보다 목표 오프셋, 속도 제한, 필터, 감쇠 설정으로 비슷한 응답을 근사하는 경우가 많다.

여기서 처음 본 식 \eqref{eq:task_impedance}의 희망 관성 $\bm{M}_d$가 사라진 것처럼 보일 수 있다.
희망 관성을 없앤 것이 아니라, 여기서는 먼저 구현하기 쉬운 가상 스프링-댐퍼 명령힘만 꺼내 본다.
희망 관성까지 맞추려면 실제 관성, 자코비안, 동역학 보상, 토크 제어 성능을 함께 다뤄야 한다.
그래서 교육용 설명에서는 먼저 강성과 감쇠가 손끝의 접촉 감각을 어떻게 바꾸는지 보여주고, 희망 관성은 뒤에서 유효 질량과 함께 조심스럽게 해석한다.

결국 로봇 끝단을 목표점에 묶인 가상의 스프링과 댐퍼처럼 느껴지게 만드는 셈이다.

부호 기준은 1차원 예에서 가장 명확하게 고정된다.
목표점이 $x_d=0$이고 실제 위치가 오른쪽으로 $x=0.01~\mathrm{m}$ 밀려 있다고 하자.
강성이 $k_d=100~\mathrm{N/m}$이면 가상 스프링이 로봇에 명령하는 힘은
\begin{equation}
  f_{\mathrm{cmd}}=-k_d(x-x_d)=-1~\mathrm{N}
\end{equation}
이다.
즉, 로봇은 왼쪽으로 $1~\mathrm{N}$만큼 되돌아오려 한다.
반대로 외부 사람이 로봇을 그 위치에 가만히 붙잡아 두려면 오른쪽으로 $1~\mathrm{N}$의 힘을 가해야 한다.
정지해서 서로 버티는 순간에는 로봇에 작용하는 명령힘과 외력이 반대 방향으로 정적 평형을 이룬다.
그래서 식 \eqref{eq:impedance_1d}에서 외력과 위치 오차의 부호는 어떤 힘을 기준으로 쓰는지에 따라 달라질 수 있지만, 물리적 의미는 같다.
평형점에서 벗어나면 가상 스프링은 그 오차를 줄이는 방향으로 작용한다.

\subsection{평형점이란 무엇인가?}

임피던스 제어에서 평형점(equilibrium point)은 가장 먼저 잡아야 할 기준이다.
평형점은 힘들이 서로 균형을 이루어 더 이상 움직이지 않는 위치이다.
다만 임피던스 제어에서는 두 종류의 평형을 구분하면 의미가 선명해진다.
평형점이라는 말을 들으면 항상 외력이 없는 장면인지, 일정한 외력이 걸린 뒤의 장면인지 함께 물어야 한다.
$x_d$는 외력이 없을 때의 기준 평형점, 즉 가상 스프링이 묶여 있는 anchor point이고, 외력이 일정하게 걸린 뒤 실제로 멈추는 위치는 뒤에서 계산할 $x_{\mathrm{ss}}$이다.

\begin{table}[!htbp]
  \centering
  \caption{임피던스 제어에서 구분해야 할 두 평형점}
  \small
  \begin{tabularx}{\linewidth}{@{}>{\raggedright\arraybackslash}p{0.22\linewidth}>{\raggedright\arraybackslash}p{0.28\linewidth}>{\raggedright\arraybackslash}X@{}}
    \toprule
    기호 & 이름 & 의미 \\
    \midrule
    $x_d$ & 무부하 평형점, anchor point & 외력이 없을 때 가상 스프링 힘이 0이 되는 기준점 \\
    $x_{\mathrm{ss}}$ & 하중이 걸린 정적 평형점 & 일정한 외력이 걸린 뒤 실제로 멈추는 위치 \\
    \bottomrule
  \end{tabularx}
\end{table}

스프링만 있는 가장 단순한 경우를 보자.
목표 위치를 $x_d$라고 하고, 스프링 힘을
\begin{equation}
  f_k=-k_d(x-x_d)
\end{equation}
라고 하자.
외력이 없다면 $x=x_d$일 때 스프링 힘이 0이 된다.
따라서 이때의 평형점은 $x_d$이다.

그래서 임피던스 제어에서 $x_d$는 단순한 목표 위치가 아니라 가상 스프링이 고정된 기준점이다.
영어로는 anchor point, equilibrium point, desired position 같은 표현이 상황에 따라 함께 쓰인다.
$x_d$는 로봇 끝단을 잡아당기는 가상 스프링의 고정점에 가깝다.

외력이 있으면 평형점은 달라진다.
정상상태에서는 속도와 가속도가 0이므로
\begin{equation}
  \dot{x}=0, \qquad \ddot{x}=0
\end{equation}
이다.
식 \eqref{eq:impedance_1d}에서 목표도 정지해 있다고 하면
\begin{equation}
  k_d(x-x_d)=f_{\mathrm{ext}}
\end{equation}
만 남는다.
따라서 외력이 일정할 때 최종 위치는
\begin{equation}
  x_{\mathrm{ss}}
  =
  x_d+\frac{f_{\mathrm{ext}}}{k_d}
  \label{eq:steady_position_under_force}
\end{equation}
가 된다.

이 식은 임피던스 제어의 성격을 매우 잘 보여준다.
임피던스 제어는 외력이 있어도 무조건 $x_d$에 머무는 제어가 아니다.
외력이 있으면 강성으로 정해지는 양만큼 밀린다.
더 정확히는 변위는 힘에 비례하고 강성에는 반비례한다.
강성 $k_d$가 크면 같은 힘을 받아도 적게 밀리고, 강성 $k_d$가 작으면 많이 밀린다.
즉, 임피던스 제어의 정상상태 위치 정확도는 강성과 외력의 관계로 결정된다.

이것이 위치 제어와의 중요한 차이이다.
강한 위치 제어는 외력이 있어도 목표 위치를 최대한 지키려 한다.
반면 임피던스 제어는 힘을 받으면 어느 정도 밀려도 되는 관계 자체를 설계한다.
그래서 사람이나 환경과 접촉하는 상황에서 접촉력 변화가 완만한 상호작용을 만들 수 있다.
다만 실제 안전성은 강성 하나로 보장되지 않고, 감쇠, 속도 제한, 힘 제한, 센서 대역폭, 환경의 강성까지 함께 좌우한다.

라플라스 변환으로 들어가기 전에 네 가지 제어 관점을 짧게 분리해 두자.
위치 제어는 어디에 있어야 하는가를 가장 중요하게 보고, 힘 제어는 얼마나 밀어야 하는가를 가장 중요하게 본다.
임피던스 제어는 둘을 절반씩 섞는 방식이 아니라, 힘을 받았을 때 얼마나 밀리고 어떤 속도로 되돌아올지를 정한다.
어드미턴스 제어는 같은 관계를 구현 반대편에서 사용해, 측정된 힘으로 목표 위치나 속도 명령을 갱신한다.

\begin{table}[!htbp]
  \centering
  \caption{처음 읽는 제어 방식 구분}
  \small
  \begin{tabularx}{\linewidth}{@{}>{\raggedright\arraybackslash}p{0.20\linewidth}>{\raggedright\arraybackslash}p{0.28\linewidth}>{\raggedright\arraybackslash}X@{}}
    \toprule
    방식 & 먼저 묻는 질문 & 손으로 밀어 보면 느껴지는 차이 \\
    \midrule
    위치 제어 & 목표 위치를 지킬 것인가? & 잘 밀리지 않으려 하고, 접촉 힘이 커질 수 있다. \\
    힘 제어 & 목표 힘을 맞출 것인가? & 정해진 힘으로 누르려 하지만, 위치는 작업 환경에 따라 달라진다. \\
    임피던스 제어 & 힘과 움직임의 관계를 어떻게 만들 것인가? & 스프링-댐퍼처럼 어느 정도 밀리고 다시 돌아온다. \\
    어드미턴스 제어 & 측정된 힘을 어떤 위치/속도 명령으로 바꿀 것인가? & 힘을 받으면 바깥 위치 제어기의 목표점이 부드럽게 움직인다. \\
    \bottomrule
  \end{tabularx}
\end{table}

\subsection{라플라스 변환과 전달함수 관점}

라플라스 변환과 전달함수는 처음에는 수식 도구처럼 보이지만, 여기서 보고 싶은 장면은 단순하다.
일정한 힘을 계단 입력처럼 갑자기 가한다.
로봇 끝단이 잠깐 흔들리거나 지나친다.
시간이 충분히 지나면 어느 위치에서 멈춘다.
전달함수는 이 전체 과정을 하나의 식으로 묶어, 최종값은 무엇이 정하고 흔들림은 무엇이 정하는지 보게 해준다.

여러 이름이 이어지므로 먼저 큰 그림부터 잡는다.
같은 질량-스프링-댐퍼 식이라도 무엇을 입력으로 보고 무엇을 출력으로 보느냐에 따라 이름이 달라진다.
이 표에서는 목표점 기준 위치 오차를 $e=x-x_d$, 그 라플라스 변환을 $E(s)$로 둔다.
영 초기조건의 전달함수 문맥에서는 오차 속도의 라플라스 변환을 $V_e(s)=sE(s)$로 쓸 수 있다.
초기 오차가 남아 있으면 엄밀히는 $V_e(s)=sE(s)-e(0)$가 되므로, 그 경우에는 초기조건이 만드는 자유응답을 따로 보아야 한다.
절대 속도는 $V_x(s)$로 구분한다.

\begin{table}[!htbp]
  \centering
  \caption{같은 동역학을 읽는 네 가지 관점}
  \small
  \setlength{\tabcolsep}{4pt}
  \begin{tabularx}{\linewidth}{@{}>{\raggedright\arraybackslash}p{0.23\linewidth}>{\raggedright\arraybackslash}p{0.20\linewidth}>{\raggedright\arraybackslash}p{0.20\linewidth}>{\raggedright\arraybackslash}X@{}}
    \toprule
    관점 & 입력 & 출력 & 쉬운 의미 \\
    \midrule
    오차 순응성 $C_e(s)$ & 외력 $F_{\mathrm{ext}}(s)$ & 위치 오차 $E(s)$ & 힘을 받으면 목표점에서 얼마나 밀리는가. 단위는 $\mathrm{m/N}$ \\
    오차 어드미턴스 $Y_e(s)$ & 외력 $F_{\mathrm{ext}}(s)$ & 오차 속도 $V_e(s)$ & 힘을 받으면 목표점 기준으로 얼마나 흘러 움직이는가. 단위는 $(\mathrm{m/s})/\mathrm{N}$ \\
    오차 임피던스 $Z_{\mathrm{err}}(s)$ & 오차 속도 $V_e(s)$ & 힘 $F_{\mathrm{ext}}(s)$ & 목표점 기준 그 속도로 밀려 움직이려면 힘이 얼마나 필요한가. 고정 목표점에서는 기계 임피던스처럼 읽을 수 있다. 단위는 $\mathrm{N}/(\mathrm{m/s})$ \\
    정적 순응성 & 일정한 힘 $F_0$ & 최종 위치 오차 $e_{\mathrm{ss}}$ & 오래 기다린 뒤 얼마나 밀려 있는가. 단위는 $\mathrm{m/N}$ \\
    \bottomrule
  \end{tabularx}
\end{table}

이번에는 같은 식을 라플라스 영역에서 보자.
라플라스 영역은 미분 연산을 $s$를 곱하는 연산으로 바꾸어, 입력과 출력의 비율을 보기 쉽게 만드는 도구이다.
전달함수(transfer function)는 이 라플라스 영역에서 입력을 넣었을 때 출력이 어떤 비율과 시간 특성으로 나오는지를 나타내는 함수이다.
예를 들어 $F_{\mathrm{ext}}(s)$를 입력으로 넣었을 때 $E(s)$가 얼마나 나오느냐를 보면, 힘이 위치 오차로 어떻게 바뀌는지 읽을 수 있다.
처음에는 대수 조작보다 분모의 세 값, $m_d$, $d_d$, $k_d$가 응답 모양을 정한다는 점만 보자.
이 절의 목표는 전달함수 계산 자체보다, 왜 최종 위치 오차는 강성이 정하고 흔들림은 질량과 감쇠가 정하는지 확인하는 것이다.
아래 전개는 목표점이 고정된 1차원 선형시불변 모델, 영 초기조건, 그리고 $m_d>0$, $d_d>0$, $k_d>0$인 안정한 경우를 가정한다.
목표점이 고정되어 있고 $e=x-x_d$라고 하면 1차원 임피던스 식은
\begin{equation}
  m_d\ddot{e}+d_d\dot{e}+k_d e=f_{\mathrm{ext}}
  \label{eq:error_impedance_1d}
\end{equation}
이다.
초기조건을 0으로 두고 라플라스 변환을 적용하면
\begin{equation}
  (m_d s^2+d_d s+k_d)E(s)=F_{\mathrm{ext}}(s)
\end{equation}
가 된다.
따라서 외력에서 위치 오차까지의 전달함수는
\begin{equation}
  \frac{E(s)}{F_{\mathrm{ext}}(s)}
  =
  \frac{1}{m_d s^2+d_d s+k_d}.
  \label{eq:force_to_position_compliance}
\end{equation}
이 전달함수는 오차 순응성(error compliance), 즉 힘을 받았을 때 목표점에서 얼마나 벗어나는지를 나타낸다.
여기서 $E(s)=\mathcal{L}\{e(t)\}$이며, 질량-스프링-댐퍼 장에서 쓴 총 에너지 $E(t)$와는 다른 기호 사용이다.
강성이 단단함의 정도라면, 순응성은 부드럽게 밀려나는 정도이다.
같은 외력에서 많이 움직이면 순응성이 큰 것이고, 거의 움직이지 않으면 순응성이 작은 것이다.
스프링 하나만 놓고 보면 위치 오차는 $e=F/k_d$이므로, 정적인 순응성은 $1/k_d$로 볼 수 있다.
단위는 $\mathrm{m/N}$이다.
강성과 정적 순응성은 스칼라 정적 모델에서 서로 역수 관계이고, 오차 순응성과 오차 어드미턴스는 출력이 위치 오차냐 오차 속도냐에 따라 $Y_e=sC_e$로 연결된다.
넓게 말하면 임피던스와 어드미턴스, 강성과 순응성은 서로 반대 방향에서 같은 현상을 바라보는 말이다.
다만 이 역수 관계는 같은 포트 변수, 스칼라 SISO LTI 모델, 그리고 역이 잘 정의되는 경우에 가장 안전하게 읽어야 한다.
위치 기준 순응성 $E/F$, 속도 기준 어드미턴스 $V/F$, 임피던스 $F/V$, 동적 강성 $F/E$는 서로 가까운 개념이지만 단위와 물리적 질문이 같지는 않다.
다만 여기서 말하는 기계적 어드미턴스라는 물리량과, 뒤에서 설명할 어드미턴스 제어라는 제어 구조는 연결되어 있지만 완전히 같은 말은 아니다.
어드미턴스는 힘에서 운동으로 가는 동역학적 관계이고, 어드미턴스 제어는 그 관계를 이용해 측정 힘으로부터 위치나 속도 명령을 만들어 내는 구현 방식이다.
단위로 점검하면 차이가 더 잘 보인다.
$C_e(s)$는 $\mathrm{m/N}$처럼 힘 1 N이 위치 오차 몇 m로 바뀌는지를 읽고, $Y_e(s)$는 $\mathrm{m/(N\,s)}$처럼 힘 1 N이 오차 속도 몇 $\mathrm{m/s}$로 바뀌는지를 읽는다.
목표점이 고정되어 있으면 $V_e(s)=V_x(s)$이므로 이 오차 어드미턴스를 일반적인 기계 어드미턴스 $Y_m(s)=V_x(s)/F_{\mathrm{ext}}(s)$처럼 읽을 수 있다.

다만 오차 순응성과 오차 어드미턴스는 단위가 다르다.
동적 오차 순응성은
\begin{equation}
  C_e(s)=\frac{E(s)}{F_{\mathrm{ext}}(s)}
\end{equation}
이고, 단위는 위치 오차/힘이다.
목표점을 원점으로 잡으면 위치 오차가 곧 위치처럼 보이므로, 고정 목표점 좌표계에서는 위치 순응성처럼 읽을 수 있다.
오차 어드미턴스는 힘에서 오차 속도까지의 관계이므로, 영 초기조건의 전달함수 문맥에서는
\begin{equation}
  Y_e(s)=\frac{V_e(s)}{F_{\mathrm{ext}}(s)}=sC_e(s)
\end{equation}
이다.
고정 목표점에서는 $V_e(s)=V_x(s)$이므로 $Y_e(s)$와 $Y_m(s)$를 같은 속도/힘 관계처럼 읽을 수 있지만, 목표점이 움직이면 $V_x(s)=V_e(s)+V_d(s)$가 되어 두 값이 달라진다.
정적 순응성 $1/k_d$는 $s\to 0$에서 위치와 힘의 최종 균형을 볼 때의 값이다.

힘에서 오차 속도까지의 반대 관계로 보면 오차 임피던스가 나온다.
속도 $V_e(s)=sE(s)$이므로
\begin{equation}
  \frac{F_{\mathrm{ext}}(s)}{V_e(s)}
  =
  m_d s+d_d+\frac{k_d}{s}.
  \label{eq:control_impedance_transfer}
\end{equation}
이것은 목표점 기준으로 본 질량-스프링-댐퍼의 오차 임피던스이다.
오차 어드미턴스를
\begin{equation}
  Y_e(s)=\frac{V_e(s)}{F_{\mathrm{ext}}(s)}
\end{equation}
라고 두면, 목표점이 고정되어 $V_e=V_x$일 때만 일반적인 기계 포트 임피던스와 같은 값으로 읽을 수 있다.
일반 기계 포트 임피던스는 실제 속도 $V_x$를 기준으로 한 $F_{\mathrm{ext}}/V_x$이다.
목표점이 움직이면 $V_e=V_x-V_d$이므로, $F_{\mathrm{ext}}V_e$는 실제 외부 포트 파워가 아니라 목표점 기준 오차 운동에 대한 계산량이다.
위치 순응성 $C_e(s)=E(s)/F_{\mathrm{ext}}(s)$의 단순 역수가 아니라, 속도 기준으로 보기 때문에
\begin{equation}
  Y_e(s)=sC_e(s)
\end{equation}
라는 $s$가 하나 더 붙는다.
질량항은 $m_d s$, 감쇠항은 $d_d$, 스프링항은 $k_d/s$로 나타난다.
여기서 $F_{\mathrm{ext}}$는 양의 오차 운동을 유지하기 위해 환경이 로봇에 가하는 힘이다.
로봇이 목표점으로 되돌아가려고 만드는 복원 명령힘은 같은 오차에 대해 반대 부호로 나타난다.
목표점이 고정되어 있을 때는 $V_x(s)$와 $V_e(s)$를 거의 같은 속도 신호처럼 읽을 수 있다.
하지만 목표점이 움직이면 반드시 상대속도 $V_e(s)=sE(s)$ 또는 $\dot{x}-\dot{x}_d$를 기준으로 보아야 한다.

정상상태와 과도응답은 장면을 나누는 말이다.
정상상태는 속도와 가속도가 0인 마지막 사진이고, 과도응답은 그 사진에 도착하기 전의 동영상이다.
마지막 사진만 보면 강성이 주로 남고, 동영상을 보면 관성, 감쇠, 강성이 모두 드러난다.

정상상태응답은 $s\to 0$ 근처의 성질로 읽을 수 있다.
앞에서 교류 회로의 정현파 정상상태는 같은 주파수로 계속 흔들리는 상태를 뜻했다.
여기서의 정상상태응답은 일정한 힘을 오래 가한 뒤 속도와 가속도가 0이 된 최종 균형을 뜻한다.
따라서 여기서 $s\to0$은 0초를 본다는 뜻이 아니라, 아주 느린 변화 또는 충분히 오래 기다린 뒤의 균형을 보는 관점이다.
예를 들어 목표점이 고정되어 있고, 일정한 외력 $F_0$가 계단 입력처럼 들어온다고 하자.
또 제어 입력이 포화되지 않고, 1축 선형시불변(linear time-invariant, LTI) 근사가 유효하며, 초기조건 때문에 생긴 과도항이 안정적으로 사라진다고 가정하자.
그러면 최종값 정리를 사용해
\begin{equation}
  e_{\mathrm{ss}}
  =
  \lim_{s\to 0}
  s
  \left[
  \frac{1}{m_d s^2+d_d s+k_d}
  \frac{F_0}{s}
  \right]
  =
  \frac{F_0}{k_d}
\end{equation}
가 된다.
부호가 먼저 헷갈린다면 처음에는 크기 관계 $|e_{\mathrm{ss}}|=|F_0|/k_d$로 읽어도 좋다.
최종값 정리는 안정한 시스템이 실제로 최종값을 가진다는 조건에서
\begin{equation}
  \lim_{t\to\infty} e(t)
  =
  \lim_{s\to 0}sE(s)
\end{equation}
로 쓸 수 있다는 정리이다.
여기서는 일정한 힘을 오래 가한 뒤의 최종 위치 오차를 보기 위해 사용한다.
이 예제에서는 $m_d>0$, $d_d>0$, $k_d>0$이므로 안정한 2차 시스템이 되고, 계단 외력에 대한 최종값이 존재한다.
더 엄밀히 말하면 이 단계에서는 $sE(s)$의 극점이 모두 왼쪽 반평면에 있어야 한다.
$k_d=0$인 경우, 불안정한 이득, 지속 진동, 포화, 시간에 따라 계속 변하는 입력에서는 이 최종값 정리를 그대로 쓰면 안 된다.
여기서 $m_d$와 $d_d$는 사라지고 $k_d$만 남는다.
왜냐하면 지금 보는 정상상태는 충분히 오래 지나 로봇이 최종 위치에서 멈춘 장면이기 때문이다.
이 장면에서는 속도와 가속도가 0에 가까워진다.
따라서 일정한 힘을 받는 최종 위치 오차는 강성으로만 결정된다.

반대로 과도응답(transient response)은 $m_d$, $d_d$, $k_d$가 모두 결정한다.
외력이 갑자기 바뀌거나 목표점이 갑자기 움직이면, 로봇은 바로 최종 위치로 순간이동하지 않는다.
가속하고, 속도를 갖고, 목표점을 지나칠 수도 있고, 진동하다가 멈출 수도 있다.
이때 질량은 움직임을 시작하고 멈추기 어렵게 만드는 관성으로 나타나고, 감쇠는 속도에 비례해 걸리는 브레이크처럼 나타나며, 강성은 목표점으로 다시 끌어당기는 복원 효과로 나타난다.
그래서 움직이는 중의 그래프 모양은 질량, 감쇠, 강성 전체가 함께 결정한다.
다만 최종 정적 오차만 따로 보면 속도와 가속도가 0이 되어 강성만 남는다.

아래의 고유진동수와 감쇠비 공식은 축들이 서로 독립인 1차원 모델, 또는 대각 행렬의 각 축을 따로 보는 경우에 해당한다.
표준 2차 시스템 형태로 쓰면
\begin{equation}
  \ddot{e}
  +
  2\zeta\omega_n\dot{e}
  +
  \omega_n^2 e
  =
  \frac{1}{m_d}f_{\mathrm{ext}}
\end{equation}
이고,
이 표준형은 목표 임피던스 오차 방정식 $m_d\ddot{e}+d_d\dot{e}+k_de=f_{\mathrm{ext}}$의 양변을 $m_d$로 나눈 것이다.
계수를 표준형과 맞추면 $\omega_n^2=k_d/m_d$, $2\zeta\omega_n=d_d/m_d$이고, 전기 시스템 장의 RLC 계수 비교와 똑같은 계산으로 $\zeta$를 분리하면
\vmark{impedance-error-normalization}%
\begin{align}
  \omega_n &= \sqrt{\frac{k_d}{m_d}}, \\
  \zeta &= \frac{d_d}{2\sqrt{m_d k_d}}.
\end{align}
여기서 $\omega_n$은 고유진동수이고, $\zeta$는 감쇠비이다.
강성 $k_d$는 최종 위치 오차와 고유진동수를 바꾸고, 감쇠비 $\zeta$는 진동이 얼마나 남는지를 바꾼다.
같은 질량에서 $k_d$가 커지면 $\omega_n$이 커져 더 빨리 움직이려 하고, 같은 강성에서 $m_d$가 커지면 $\omega_n$이 작아져 더 느리게 움직이려 한다.
또 같은 감쇠계수 $d_d$라도 $m_d$와 $k_d$가 바뀌면 $\zeta$가 달라지므로, 같은 댐퍼 숫자가 항상 같은 흔들림을 뜻하지는 않는다.
단, $\omega_n$은 무감쇠 기준 고유진동수이다.
부족감쇠 응답에서 실제 피크 사이 간격으로 보이는 진동은 $\omega_d=\omega_n\sqrt{1-\zeta^2}$에 더 가깝고, $\zeta\ge1$이면 반복해서 흔들리는 주기가 나타나지 않는다.
또 로봇의 작업공간 모델이 행렬로 결합되어 있으면, 각 축을 따로 떼어 $\omega_n$과 $\zeta$ 하나씩 붙이는 해석이 정확하지 않을 수 있다.
그 경우 무감쇠 모드는 $\bm{K}_d\bm{\phi}_i=\omega_i^2\bm{M}_d\bm{\phi}_i$의 일반화 고유값으로 읽고, 감쇠가 있으면 실제 pole은 $\det(s^2\bm{M}_d+s\bm{D}_d+\bm{K}_d)=0$ 또는 대응 state-space 행렬의 고유값으로 정해진다.

여기서 얻고 싶은 것은 라플라스 계산 자체가 아니다.
오래 힘을 가한 뒤의 최종 오차는 주로 강성이 정하고, 목표점으로 가는 동안의 흔들림은 고유진동수와 감쇠비가 정한다는 감각이다.
다만 실제 Lab04 데모에는 포화, 접촉 경계, DLS 목표 오프셋, position actuator 인터페이스가 함께 들어간다.
DLS는 damped least squares의 약자로, 특이점 근처에서 역기구학 계산이 지나치게 예민해지지 않도록 감쇠를 넣어 푸는 방법이다.
DLS 목표 오프셋은 자코비안 역문제를 안정적으로 풀기 위해 목표 위치 보정을 조금씩 관절 목표로 바꾸는 값이고, position actuator는 토크가 아니라 목표 위치를 받는 구동기 인터페이스이다.
따라서 이 식들은 실행 결과를 보증하는 공식이라기보다, 플롯을 읽기 위한 첫 기준선으로 쓰는 편이 좋다.
이 구분을 잡고 나면 다음 절의 왜 강성과 감쇠비를 주로 만지는가라는 질문이 자연스럽게 이어진다.

\subsection{왜 강성과 감쇠비를 주로 설정하는가?}

이론적으로는 임피던스 제어에서 $m_d$, $d_d$, $k_d$를 모두 설계할 수 있다.
식 \eqref{eq:task_impedance}에도 분명히 희망 관성, 희망 감쇠, 희망 강성이 모두 들어 있다.
그런데 실제 로봇 제어기나 교육용 데모에서는 사용자가 직접 만지는 파라미터가 보통 강성 $k_d$와 감쇠비 $\zeta$인 경우가 많다.
질량이나 감쇠계수 $d_d$를 직접 입력하지 않는 이유는 여러 가지가 있다.
이제 앞의 기준선을 가지고 실제 튜닝 화면을 떠올려 보자.
사용자가 매번 질량, 감쇠계수, 강성을 모두 직접 고르는 방식은 생각보다 불편하다.
그래서 많은 데모와 제어기는 강성과 감쇠비를 앞에 내세운다.

특히 질량은 이름은 하나지만, 실제로는 여러 뜻이 겹쳐 쓰인다.
이 구분을 놓치면 질량을 $1~\mathrm{kg}$으로 둔다는 말이 실제 로봇이 1 kg짜리 물체처럼 움직인다는 뜻으로 오해되기 쉽다.
\begin{table}[!htbp]
  \centering
  \caption{임피던스 제어에서 헷갈리기 쉬운 질량 기준}
  \small
  \setlength{\tabcolsep}{4pt}
  \begin{tabularx}{\linewidth}{@{}>{\raggedright\arraybackslash}p{0.25\linewidth}>{\raggedright\arraybackslash}p{0.27\linewidth}>{\raggedright\arraybackslash}X@{}}
    \toprule
    이름 & 의미 & 왜 구분해야 하는가 \\
    \midrule
    실제 로봇 질량/관성 & 링크, 모터, 감속기 등 로봇 하드웨어가 가진 물리적 관성 & 사용자가 마음대로 바꿀 수 없고 자세에 따라 동역학 효과가 달라진다 \\
    끝단 유효질량 & 특정 자세와 방향에서 손끝을 밀 때 느껴지는 겉보기 질량 & 같은 로봇도 방향과 자세에 따라 무겁거나 가볍게 느껴질 수 있다 \\
    희망 관성 $m_d$ 또는 $\bm{M}_d$ & 제어기가 목표로 삼는 가상 질량-스프링-댐퍼 모델의 관성 & 완전히 구현하려면 동역학 보상과 토크 제어 성능이 필요하다 \\
    교육용 명목 질량 $m_{\mathrm{nom}}$ & 감쇠계수 계산을 단순화하려고 내부에서 정해 두는 기준 질량 & $d_d=2\zeta\sqrt{m_{\mathrm{nom}}k_d}$ 같은 계산에는 편리하지만 실제 끝단 감쇠비를 보증하지는 않는다 \\
    \bottomrule
  \end{tabularx}
\end{table}

첫째, 강성은 사용자가 가장 직관적으로 느끼는 파라미터이다.
강성은 얼마나 딱딱한지, 같은 힘을 받았을 때 얼마나 밀리는지를 정한다.
정상상태에서
\begin{equation}
  e_{\mathrm{ss}}=\frac{F_0}{k_d},
  \qquad
  |e_{\mathrm{ss}}|=\frac{|F_0|}{k_d}
\end{equation}
이므로, 강성은 접촉 중 위치 오차와 접촉력의 크기에 직접 연결된다.
예를 들어 $k_d=1000~\mathrm{N/m}$이면 $10~\mathrm{N}$의 힘을 받을 때 약 $1~\mathrm{cm}$ 밀린다.
$k_d=100~\mathrm{N/m}$이면 같은 힘에서 약 $10~\mathrm{cm}$ 밀린다.
이 차이는 사람이 손으로 만져도 바로 느낄 수 있다.

감쇠비는 흔들림을 읽는 가장 간단한 눈금이다.
정확한 값은 시스템과 입력에 따라 달라지지만, 처음 튜닝할 때는 다음 표처럼 응답 모양을 떠올리면 출발점이 생긴다.
\begin{table}[!htbp]
  \centering
  \caption{감쇠비를 응답 모양으로 읽는 법}
  \small
  \setlength{\tabcolsep}{4pt}
  \begin{tabularx}{\linewidth}{@{}>{\raggedright\arraybackslash}p{0.18\linewidth}>{\raggedright\arraybackslash}p{0.30\linewidth}>{\raggedright\arraybackslash}X@{}}
    \toprule
    감쇠비 & 응답 모양 & 튜닝 감각 \\
    \midrule
    $\zeta<1$ & 부족감쇠, 흔들림과 오버슈트 가능 & 빠르지만 튈 수 있다 \\
    $\zeta\approx0.7$ & 비교적 빠르고 오버슈트가 제한됨 & 실험 출발값으로 자주 쓰기 좋다 \\
    $\zeta\approx1$ & 임계감쇠 근처, 오버슈트가 작음 & 접촉 충격과 진동을 줄이고 싶을 때 유용하다 \\
    $\zeta>1$ & 과감쇠, 반복 진동은 적지만 느림 & 안전하게 보일 수 있으나 둔하고 effort가 커질 수 있다 \\
    \bottomrule
  \end{tabularx}
\end{table}

둘째, 감쇠계수 $d_d$ 자체보다 감쇠비 $\zeta$가 더 의미 있는 설계 변수이다.
같은 $d_d=20~\mathrm{N\,s/m}$라도 질량과 강성이 달라지면 응답은 완전히 달라진다.
가벼운 시스템에서는 매우 큰 감쇠처럼 느껴질 수 있고, 무거운 시스템에서는 부족한 감쇠일 수 있다.
반면 감쇠비는 임계감쇠에 대한 상대적인 비율이다.
\begin{equation}
  d_d=2\zeta\sqrt{m_d k_d}
  \label{eq:damping_from_ratio}
\end{equation}
로 계산하면, 강성이나 질량이 바뀌어도 진동을 어느 정도 줄일 것인가라는 의도가 비교적 일관되게 유지된다.
그래서 사용자는 감쇠계수 37.2 같은 숫자를 외우기보다 감쇠비 0.7, 감쇠비 1.0처럼 응답 모양을 기준으로 설정하는 편이 쉽다.
다만 식 \eqref{eq:damping_from_ratio}은 고정된 양의 스칼라 질량을 갖는 1축 2차 시스템이나, 축을 독립적으로 볼 수 있는 대각 모델에서 가장 깔끔하게 해석된다.
행렬 임피던스에서는 출발값으로 쓰고, 필요한 경우 모드별 감쇠나 행렬 기반 설계를 따로 확인해야 한다.
이제부터는 실제 로봇에서 어떤 근사가 들어가는지 보는 구간이다.
교육용 데모에서는 내부적으로 몇 가지 단순 정책을 쓸 수 있다.
가장 단순하게는 각 축의 명목 질량을 $1~\mathrm{kg}$처럼 고정해 두고 $d_d$를 계산한다.
조금 더 로봇답게 만들려면 현재 자세에서 끝단 유효질량을 추정해 그 값으로 $d_d$를 계산할 수 있다.
또는 하드웨어 인터페이스가 위치 제어형이라면, 실제 관성 구현은 포기하고 목표 위치 필터나 속도 감쇠로 비슷한 응답만 만들 수도 있다.
어떤 경우든 내부에서 사용한 질량이 실제 끝단 유효질량과 다르면, 설정한 감쇠비는 정확한 물리 감쇠비가 아니라 교육용 또는 제어기 내부의 근사값으로 읽어야 한다.
그래서 화면에는 보통 강성 $k_d$와 감쇠비 $\zeta$가 먼저 보이고, 감쇠계수 $d_d$는 선택한 $\zeta$와 내부 질량 기준으로 계산되어 제어식 안에 들어간다.

셋째, 로봇의 실제 관성은 사용자가 마음대로 바꾸기 어렵다.
로봇 팔의 끝단에서 느껴지는 관성은 자세에 따라 달라진다.
관절이 접혀 있을 때와 뻗어 있을 때, 같은 끝단 방향이라도 유효 질량이 달라질 수 있다.
이를 작업공간 관성, 또는 겉보기 관성(apparent inertia)이라고 부른다.
말 그대로 로봇 끝단에서 느껴지는 관성이다.
같은 로봇이라도 자세와 힘 방향이 달라지면 손끝이 무겁게 느껴지는 방향과 가볍게 느껴지는 방향이 달라질 수 있다.
에너지로 보면 자세가 관성을 바꾸는 이유가 먼저 선명해진다.
손끝을 같은 속도 $\dot{\bm{x}}$로 움직이더라도, 현재 자세에 따라 필요한 관절 속도 $\dot{\bm{q}}$의 조합이 달라진다.
관절공간 운동 에너지는
\begin{equation}
  T_q=\frac{1}{2}\dot{\bm{q}}^{\mathsf{T}}\bm{M}_q(\bm{q})\dot{\bm{q}}
\end{equation}
이므로, 같은 손끝 속도라도 어떤 자세에서는 여러 관절을 많이 움직여 큰 운동 에너지가 필요하고, 어떤 자세에서는 상대적으로 작은 운동 에너지로 같은 손끝 속도를 만들 수 있다.
사람이 로봇 끝단을 밀 때 어떤 방향은 무겁고 어떤 방향은 가볍게 느끼는 이유가 바로 여기에 있다.
조금 엄밀히 말하면, 주어진 손끝 속도 방향의 운동에너지 기준 질량은 작업공간 관성 $\bm{\Lambda}=(\bm{J}_v\bm{M}_q^{-1}\bm{J}_v^{\mathsf{T}})^{-1}$가 존재할 때 $\bm{n}^{\mathsf{T}}\bm{\Lambda}\bm{n}$처럼 읽는다.
반면 아래에 쓰는 흔한 스칼라 식은 같은 방향으로 힘을 가했을 때 그 방향 가속도가 얼마나 나오는지를 보는 force-to-acceleration 겉보기 질량이다.
여기서부터는 고급 해석을 위한 미리보기이다.
관절공간 관성행렬이 양의 정부호라는 말은 어떤 관절 방향으로도 관성 에너지가 음수가 되지 않는다는 뜻이고, 자코비안 랭크가 충분하다는 말은 그 방향의 손끝 운동을 관절 움직임으로 만들 수 있다는 뜻이다.
자유공간에 있는 조작기를 같은 좌표계에서 보고, 관절공간 관성행렬 $\bm{M}_q$가 양의 정부호이며, 병진 자코비안 $\bm{J}_v$가 해당 방향에서 충분한 랭크를 갖는다고 하자.
$\bm{J}_v\bm{M}_q^{-1}\bm{J}_v^{\mathsf{T}}$라는 조합이 어디서 나오는지도 한 단계씩 따라가 볼 수 있다.
정지 상태 근처에서 속도 관련 항과 중력 보상을 잠시 무시하면 관절공간 운동방정식은 $\bm{M}_q\ddot{\bm{q}}=\bm{\tau}$로 줄어든다.
손끝에 힘 $\bm{f}$를 가하는 것은 로봇공학 기초 장에서 일률로 유도한 자코비안 전치 관계에 따라 관절 effort $\bm{\tau}=\bm{J}_v^{\mathsf{T}}\bm{f}$에 해당하므로, 관절 가속도는 $\ddot{\bm{q}}=\bm{M}_q^{-1}\bm{J}_v^{\mathsf{T}}\bm{f}$이다.
손끝 가속도는 가속도 관계 $\ddot{\bm{x}}=\bm{J}_v\ddot{\bm{q}}+\dot{\bm{J}}_v\dot{\bm{q}}$에서 정지 상태 $\dot{\bm{q}}=\bm{0}$이라 둘째 항이 사라지므로
\[
  \ddot{\bm{x}}
  =
  \bm{J}_v\ddot{\bm{q}}
  =
  \left(\bm{J}_v\bm{M}_q^{-1}\bm{J}_v^{\mathsf{T}}\right)\bm{f}
\vmark{lambda-force-to-acceleration-derivation}%
\]
가 된다.
$f=ma$에서 질량이 가속도를 힘으로 바꾸는 표라면, $\bm{J}_v\bm{M}_q^{-1}\bm{J}_v^{\mathsf{T}}$는 반대로 힘을 가속도로 바꾸는 표이므로 질량의 역수 역할을 한다.
그래서 그 역행렬 $\bm{\Lambda}=(\bm{J}_v\bm{M}_q^{-1}\bm{J}_v^{\mathsf{T}})^{-1}$가 작업공간에서 질량 행렬 역할을 맡는다.
힘을 가하는 방향을 단위벡터 $\bm{n}$으로 잡으면, 한 가지 흔한 표기에서는 그 방향의 끝단 유효질량을
\begin{equation}
  m_{\mathrm{eff},n}
  =
  \frac{1}
  {\bm{n}^{\mathsf{T}}\bm{J}_v\bm{M}_q^{-1}\bm{J}_v^{\mathsf{T}}\bm{n}}
\end{equation}
처럼 읽을 수 있다.
즉 $\bm{f}=f_n\bm{n}$을 가했을 때 방향 가속도 $a_n=f_n\bm{n}^{\mathsf{T}}\bm{J}_v\bm{M}_q^{-1}\bm{J}_v^{\mathsf{T}}\bm{n}$가 나온다고 보고, $f_n=m_{\mathrm{eff},n}a_n$ 꼴로 다시 쓴 값이다.
장난감 숫자로 확인해 보자.
로봇공학 기초 장의 2링크 예제 자세($l_1=l_2=1~\mathrm{m}$, $q_1=0$, $q_2=\pi/2$)에서 병진 자코비안은 첫 행 $[-1\ \ -1]$, 둘째 행 $[1\ \ 0]$이었다.
두 관절의 관성을 각각 $1~\mathrm{kg\,m^2}$로 두면($\bm{M}_q=\bm{I}$) $\bm{J}_v\bm{M}_q^{-1}\bm{J}_v^{\mathsf{T}}$는 첫 행 $[2\ \ -1]$, 둘째 행 $[-1\ \ 1]$인 행렬이 된다.
$x$방향 $\bm{n}=[1\ \ 0]^{\mathsf{T}}$에서는 $\bm{n}^{\mathsf{T}}(\bm{J}_v\bm{M}_q^{-1}\bm{J}_v^{\mathsf{T}})\bm{n}=2$이므로 $m_{\mathrm{eff}}=1/2$, $y$방향 $\bm{n}=[0\ \ 1]^{\mathsf{T}}$에서는 $1$이므로 $m_{\mathrm{eff}}=1$이다.
같은 자세라도 손끝이 $x$방향으로는 가볍게, $y$방향으로는 두 배 무겁게 느껴진다는 뜻이다.
\vmark{lambda-directional-mass-numeric-check}%
이 식은 같은 로봇이라도 자세와 방향에 따라 느껴지는 질량이 달라진다는 사실을 잘 보여준다.
특이점 근처나 수치 조건이 나쁜 자세에서는 이 값이 매우 커지거나 계산이 불안정해질 수 있으므로, 실제 구현에서는 근사값, 제한값, 또는 명목 질량을 쓰는 경우가 많다.
정확한 희망 관성 $\bm{M}_d$를 구현하려면 로봇의 동역학 모델, 자코비안, 관성행렬, 토크 제어 대역폭이 모두 충분히 좋아야 한다.
자코비안은 앞의 로봇공학 기초 장에서 관절공간과 작업공간을 연결하는 도구로 설명했다.
많은 실용 제어기에서도 실제 로봇의 관성은 내부 시스템에 맡기고 사용자가 강성과 감쇠를 조절하게 만드는 경우가 많다.
특히 위치 제어형 인터페이스만 제공되는 로봇에서는 사용자가 실제 끝단 관성을 직접 바꾸는 것이 거의 불가능하다.
이 경우 제어기는 희망 관성을 완벽히 구현한다기보다, 목표 위치를 얼마나 부드럽게 움직일지, 접촉 방향에서 얼마나 밀릴지, 속도에 얼마나 저항할지를 조절하는 방식으로 임피던스적 거동을 근사한다.
그래서 초반에는 질량 슬라이더를 노출하기보다, 먼저 강성과 감쇠비가 응답을 어떻게 바꾸는지 보여주는 편이 더 직관적이다.

넷째, 감쇠계수는 강성과 질량에 종속해서 잡는 것이 안전한 경우가 많다.
강성을 올리면 시스템은 더 빠르게 흔들리려 한다.
그런데 감쇠계수를 그대로 두면 감쇠비가 낮아져 진동이 커질 수 있다.
따라서 강성을 바꿀 때마다 식 \eqref{eq:damping_from_ratio}으로 감쇠계수를 다시 계산하는 방식이 더 자연스럽다.
사용자에게는 $k_d$와 $\zeta$를 보여주고, 내부적으로는 $d_d$를 계산하는 것이다.

실제 튜닝에서 가장 먼저 드러내기 좋은 변수는 강성과 감쇠비이다.
강성은 손끝이 얼마나 단단하게 버티는지를 정하고, 감쇠비는 그 단단함으로 움직일 때 흔들림이 얼마나 남는지를 정한다.
질량은 중요하지 않아서 빠지는 것이 아니라, 실제 로봇의 관성과 깊게 얽혀 있어 초반 튜닝 변수로 드러내기 어렵다.

\subsection{왜 스프링만으로 단순화할 수 있는가?}

이제 움직이는 동안의 이야기는 잠시 접어 두고, 로봇 손끝이 벽에 닿아 멈춘 장면만 보자.
정지 목표점과 일정한 외력 아래에서 시스템이 안정적으로 수렴했다면, 속도와 가속도는 0에 가깝다.
그 순간에는 질량항과 감쇠항이 사라지고, 힘과 변위를 잇는 스프링 항만 남는다.
또한 제어 입력이 포화되지 않고, 목표점이 계속 움직이거나 주기적으로 흔들리는 상황이 아니라고 가정한다.
그래서 정적 접촉을 처음 이해할 때는 스프링 하나로 시작해도 충분히 좋은 직관을 얻을 수 있다.
식
\begin{equation}
  m_d\ddot{e}+d_d\dot{e}+k_d e=f_{\mathrm{ext}}
\end{equation}
에서 정상상태를 가정하면
\begin{equation}
  \ddot{e}=0, \qquad \dot{e}=0
\end{equation}
이므로
\begin{equation}
  k_d e=f_{\mathrm{ext}}
\end{equation}
만 남는다.
따라서 완전히 정지한 최종 접촉력과 최종 변위를 설명할 때는 스프링 모델만으로도 핵심을 잘 보여줄 수 있다.
천천히 움직이는 준정적(quasi-static) 접촉으로 넓혀 말하려면, 감쇠항이 스프링항에 비해 충분히 작다는 조건도 함께 붙어야 한다.
\begin{equation}
  d_d|\dot{e}| \ll k_d|e|
\end{equation}
가 대략 성립할 때는 움직임이 조금 남아 있어도 스프링 직관이 주로 보인다.

예를 들어 로봇 끝단이 벽을 천천히 누르고 있다고 하자.
속도는 거의 0이고, 가속도도 거의 0이다.
이때 벽과 로봇 사이의 힘은 대부분 목표점에서 얼마나 밀렸는지와 강성이 얼마인지로 결정된다.
\begin{equation}
  |f_{\mathrm{contact}}| \approx k_d |e|.
\end{equation}
그래서 임피던스 제어를 처음 설명할 때는 가상 스프링 하나만으로도 많은 직관을 만들 수 있다.
목표점을 벽 뒤쪽으로 조금 밀어 넣으면 로봇은 벽을 누르려 하고, 강성을 낮추면 부드럽게 밀리고, 강성을 높이면 단단하게 버틴다.

하지만 스프링만으로 모든 것을 설명할 수 있는 것은 아니다.
스프링은 에너지를 저장했다가 다시 돌려준다.
그래서 스프링만 있으면 목표점 근처에서 진동할 수 있다.
댐퍼는 이 진동 에너지를 줄이는 역할을 한다.
또 질량은 로봇이 얼마나 빠르게 가속하고, 충돌 순간 얼마나 큰 동역학적 힘이 생기는지를 결정한다.
따라서 스프링만 남는다는 말은 정상상태 힘-위치 관계를 설명할 때 맞는 말이고, 과도응답이나 충돌, 빠른 움직임까지 설명하려면 질량과 감쇠가 반드시 필요하다.

그렇다면 왜 댐퍼만으로 접근하지 않을까?
댐퍼 힘은
\begin{equation}
  f_d=-d_d\dot{x}
\end{equation}
이다.
속도가 0이면 댐퍼 힘도 0이다.
즉, 댐퍼는 정지해 있는 로봇을 목표점으로 끌고 가지 못하고, 일정한 접촉력을 유지하지도 못한다.
댐퍼는 어디에 있어야 하는지를 알려주는 요소가 아니라, 움직이는 동안 에너지를 줄이는 요소이다.

질량만으로 접근하지 않는 이유도 비슷하다.
질량 힘은
\begin{equation}
  f_m=m_d\ddot{x}
\end{equation}
이다.
가속도가 0이면 질량항도 0이다.
질량은 속도 변화를 어렵게 만들지만, 어느 위치가 목표인지 알려주지 않는다.
따라서 로봇에게 목표 위치 주변의 상호작용을 만들고 싶다면 정적 위치 기준은 스프링이 정한다.
하지만 과도응답, 접촉 안정성, 충격에 대한 반응은 감쇠와 관성이 함께 정한다.

\subsection{같은 힘이라도 강성과 변위가 다르면 왜 다른가?}

부호를 잠시 내려놓고 보면, 스프링을 어떤 변위에 유지하는 외력의 크기 또는 복원력의 크기는
\begin{equation}
  |f|=k|x|
\end{equation}
로 쓸 수 있다.
여기서 $x$는 절대 위치 좌표가 아니라 스프링이 기준점에서 늘어나거나 눌린 변위이다.
로봇 임피던스 식의 표기와 맞추면 $x$ 대신 위치 오차 $e$라고 써도 된다.
부호 있는 복원력으로 쓰면 $f_k=-kx$이고, 양의 방향으로 $x$만큼 밀린 물체를 그 자리에 붙잡아 두는 외력은 $+kx$이다.
이때 처음 공부할 때 흔히 떠오르는 질문이 있다.
힘 $f$가 같다면 같은 상황으로 보아도 될까?
예를 들어
\begin{align}
  k&=1000~\mathrm{N/m}, \quad x=0.01~\mathrm{m}, \\
  k&=100~\mathrm{N/m}, \quad x=0.10~\mathrm{m}
\end{align}
은 둘 다 힘의 크기가 $f=10~\mathrm{N}$이다.
그 순간의 힘 크기만 보면 같다.
그 순간의 힘 크기만 같을 뿐, 에너지와 주변의 단단함은 다르다.

첫째, 저장된 퍼텐셜 에너지가 다르다.
스프링에 저장된 에너지는
\begin{equation}
  U=\frac{1}{2}kx^2
\end{equation}
이다.
$f=kx$를 이용하면 같은 식을
\begin{equation}
  U=\frac{f^2}{2k}
  =
  \frac{1}{2}fx
\end{equation}
로도 쓸 수 있다.
여기서 $\frac{1}{2}$가 붙는 이유는 스프링 힘이 처음부터 $f$로 고정되어 있는 것이 아니라, 변위가 0일 때 0에서 시작해 최종 변위 $x$에서 $f$까지 선형으로 증가하기 때문이다.
따라서 일을 계산할 때 평균 힘은 $f/2$이고, 힘-변위 그래프로 보면 넓이가 직사각형 $fx$가 아니라 삼각형 $\frac{1}{2}fx$가 된다.
같은 힘 크기 $f$에서 강성 $k$가 작고 변위 $x$가 크면 저장 에너지는 더 크다.
위 예시에서
\begin{align}
  U_1&=\frac{1}{2}\cdot 1000\cdot 0.01^2=0.05~\mathrm{J}, \\
  U_2&=\frac{1}{2}\cdot 100\cdot 0.10^2=0.5~\mathrm{J}
\end{align}
이다.
둘 다 그 순간 힘의 크기는 $10~\mathrm{N}$이지만, 두 번째 경우가 에너지를 10배 더 많이 저장하고 있다.
따라서 손을 놓았을 때 되돌아오며 낼 수 있는 운동도 달라진다.
앞의 경우는 단단해서 조금만 움직여도 큰 힘이 생긴 것이고, 뒤의 경우는 부드러워서 많이 밀린 뒤에야 같은 힘이 생긴 것이다.
그래서 접촉을 볼 때는 힘의 크기만 보지 말고, 그 힘이 어떤 강성과 어떤 위치 오차에서 나온 힘인지 함께 봐야 한다.

둘째, 힘의 민감도가 다르다.
스프링에서 유지 외력의 위치 기울기 또는 힘 크기의 위치 기울기는
\begin{equation}
  \frac{\dd f}{\dd x}=k
\end{equation}
이다.
부호 있는 복원력 $f_k=-kx$의 기울기는 $-k$이지만, 여기서 비교하는 것은 힘의 크기가 위치 변화에 얼마나 민감하게 늘어나는지이다.
강성이 큰 스프링은 조금만 더 밀어도 힘이 크게 증가한다.
강성이 작은 스프링은 많이 움직여도 힘이 천천히 증가한다.
따라서 같은 $10~\mathrm{N}$ 지점에 있더라도, 그 주변에서 느껴지는 단단함은 완전히 다르다.

셋째, 접촉 안정성과 위치 오차가 달라진다.
강성이 큰 시스템은 목표 위치를 잘 지키지만, 접촉 순간 힘이 급격히 커질 수 있다.
강성이 작은 시스템은 위치 오차가 커지지만, 접촉력이 부드럽게 증가한다.
즉, $f=kx$에서 $f$만 같은 것은 한 순간의 결과가 같다는 뜻일 뿐이고, 에너지, 민감도, 안전성, 조작감은 $k$와 $x$의 조합에 따라 달라진다.

\subsection{임피던스 제어, 관절 임피던스 제어, 직교 임피던스 제어}

먼저 체감 차이부터 잡자.
관절 임피던스는 각 관절에 가상의 스프링과 댐퍼를 다는 생각에 가깝다.
직교 임피던스는 손끝의 $x$, $y$, $z$ 방향에 가상의 스프링과 댐퍼를 다는 생각에 가깝다.
그래서 전자는 모터가 있는 관절 언어로 자연스럽고, 후자는 사람이 보는 작업 방향 언어로 자연스럽다.

절~\ref{sec:robotics-foundations}에서 본 것처럼, 로봇 손끝의 운동은 관절공간과 작업공간 사이의 번역을 거쳐야 한다.
지금부터의 핵심은 새 물리를 추가하는 것이 아니라, 손끝에서 설계한 스프링-댐퍼 감각이 실제 모터 관절 명령으로 번역되어야 한다는 점이다.
개념은 좌표 변환에서 출발한다.
다만 실제 로봇에서는 그 변환만으로 끝나지 않는다.
필요한 토크를 낼 수 있는지, 제어 주기가 충분히 빠른지, 특이점이나 포화가 응답을 망치지 않는지, 여유 자유도 선택이 이상한 자세를 만들지 않는지도 함께 확인해야 한다.
앞에서는 손끝에서 원하는 힘과 움직임을 말했지만, 실제 로봇은 관절 모터를 통해 움직인다.
따라서 손끝에서 정의한 힘과 움직임은 실제 모터가 이해하는 관절각, 관절속도, 관절토크로 변환되어야 한다.
앞 절의 표기처럼 더 일반적으로는 \(\bm{\tau}\)를 관절좌표와 짝을 이루는 일반화 관절 effort로 읽고, 이 절에서는 회전 관절 중심의 토크 제어 구현을 설명하므로 그 effort를 관절 토크라고 부른다.
\vmark{scope-s6-general-tau-effort}%
\vmark{scope-s6-revolute-torque-context}%
\vmark{scope-s6-torque-ladder-preserved}%
이 변환의 중심에 자코비안이 있다.

임피던스 제어는 큰 범주이고, 방금 본 것처럼 스프링-댐퍼를 어느 좌표계에 다느냐가 관절공간 임피던스와 직교공간 임피던스를 가른다.
직교공간 임피던스는 보통 Cartesian impedance 또는 task-space impedance라고 부른다.

절~\ref{sec:robotics-foundations}에서 자세한 유도를 이미 했으므로, 여기서는 임피던스 구현에 필요한 세 줄만 다시 붙잡자.
첫째, $\bm{q}$는 관절공간 언어이고, 이 절의 $\bm{x}\in\R^3$는 손끝의 병진 위치만 뜻한다.
손끝 자세까지 포함한 포즈를 다루려면 회전 오차, 렌치, 기하학적 자코비안이 따로 필요하다.
둘째, 병진 자코비안 $\bm{J}_v$는 관절 속도를 손끝 병진 속도로 옮긴다.
셋째, 자코비안 전치 $\bm{J}_v^{\mathsf{T}}$는 손끝에 만들려는 병진 힘을 관절 effort 쪽으로 되돌린다.
\vmark{scope-s6-jtf-effort-back}%
이 절의 2자유도 팔과 Panda처럼 회전 관절만 보면 그 effort는 토크로 읽으면 된다.
\vmark{scope-s6-read-as-torque}%
이 세 줄이 관절 임피던스와 직교 임피던스를 나누는 실질적인 기준이다.

\begin{table}[!htbp]
  \centering
  \caption{Section 6에서 다시 쓰는 자코비안 복습}
  \label{tab:jacobian-recap-for-impedance}
  \small
  \setlength{\tabcolsep}{4pt}
  \renewcommand{\arraystretch}{1.12}
  \begin{tabularx}{\linewidth}{@{}p{0.22\linewidth}p{0.28\linewidth}X@{}}
    \toprule
    필요한 관계 & 식 & 이 절에서의 역할 \\
    \midrule
    속도 변환 & $\dot{\bm{x}}=\bm{J}_v(\bm{q})\dot{\bm{q}}$ & 관절이 움직일 때 손끝이 어느 방향으로 얼마나 빨리 움직이는지 읽는다. \\
    힘--관절 effort 변환 & $\bm{\tau}=\bm{J}_v(\bm{q})^{\mathsf{T}}\bm{f}$ & 손끝에 만들고 싶은 병진 힘을 각 관절의 일반화 effort 기여로 나눈다. 회전 관절에서는 이 값이 토크이다. \\
\vmark{scope-s6-effort-table-label}%
\vmark{scope-s6-generalized-contribution}%
\vmark{scope-s6-revolute-table-torque}%
    좌표계 확인 & $\bm{f}$, $\dot{\bm{x}}$, $\bm{J}_v$의 기준 좌표계 일치 & 힘 센서값이나 목표 힘이 tool frame에 있으면 먼저 같은 Cartesian frame으로 회전 변환한다. \\
    \bottomrule
  \end{tabularx}
\end{table}

전치가 나오는 이유를 한 줄만 다시 적으면, 같은 순간의 손끝 일률과 관절 일률을 같게 두기 때문이다.
즉 절~\ref{sec:robotics-foundations}의 식~\eqref{eq:jacobian-transpose-foundation}은
\[
  \bm{f}^{\mathsf{T}}\dot{\bm{x}}
  =
  \bm{\tau}^{\mathsf{T}}\dot{\bm{q}}
\]
라는 가상일 관점에서 나온 결과이고, 여기서는 그 결과를 직교 임피던스 명령에 적용한다.

관절공간 임피던스(joint-space impedance)는 관절각 $\bm{q}$에 대해 임피던스를 건다.
목표 관절각을 $\bm{q}_d$, 목표 관절속도를 $\dot{\bm{q}}_d$라고 하면 관절공간 스프링-댐퍼 제어는 보통
\begin{equation}
  \bm{\tau}_{\mathrm{cmd}}
  =
  -\bm{K}_q(\bm{q}-\bm{q}_d)
  -
  \bm{D}_q(\dot{\bm{q}}-\dot{\bm{q}}_d)
  +
  \bm{\tau}_{\mathrm{bias}}
  \label{eq:joint_impedance}
\end{equation}
이다.
여기서 $\bm{\tau}_{\mathrm{bias}}$는 중력, 코리올리, 원심력 등 모델 기반 보상항을 의미한다.
목표 자세를 가만히 유지하는 가장 단순한 경우에는 $\dot{\bm{q}}_d=\bm{0}$이므로 감쇠항이 $-\bm{D}_q\dot{\bm{q}}$처럼 보인다.
하지만 관절 궤적을 따라가야 하는 경우에는 절대 속도 $\dot{\bm{q}}$가 아니라 속도 오차 $\dot{\bm{q}}-\dot{\bm{q}}_d$를 감쇠해야 한다.
관절 임피던스는 직교 임피던스보다 구현이 비교적 단순한 편이다.
다만 안정성은 이득, 포화, 샘플링 시간, 모델 보상에 좌우된다.
하지만 사용자가 느끼는 끝단의 $x$, $y$, $z$ 방향 강성은 로봇 자세에 따라 달라진다.
어떤 자세에서는 손끝이 위아래로는 부드럽고 좌우로는 단단하게 느껴질 수 있다.
관절에 스프링을 걸었기 때문에 끝단에서 보이는 스프링의 방향과 크기가 자코비안에 의해 변하기 때문이다.
예를 들어 2링크 팔의 팔꿈치가 접혀 있을 때와 거의 일직선으로 펴져 있을 때를 비교해보자.
관절 스프링 값은 같아도, 손끝에서 특정 방향으로 밀었을 때 각 관절이 움직일 수 있는 레버암과 방향이 달라진다.
그래서 같은 관절 강성이라도 손끝에서 느끼는 작업공간 강성은 자세에 따라 달라진다.

직교 임피던스(Cartesian impedance)는 로봇 끝단의 병진 위치 $\bm{x}$에 대해 임피던스를 건다.
자세까지 포함한 6차원 pose impedance를 다루려면 별도의 자세 오차와 렌치 표기가 필요하므로, 여기서는 먼저 병진 위치에 집중한다.
예를 들어 끝단 위치 오차에 대해
\begin{equation}
  \bm{f}_{\mathrm{cmd}}
  =
  -\bm{K}_d(\bm{x}-\bm{x}_d)
  -
  \bm{D}_d(\dot{\bm{x}}-\dot{\bm{x}}_d)
\end{equation}
를 계산한다.
목표점이 정지해 있다면 $\dot{\bm{x}}_d=\bm{0}$이므로 감쇠항은 보통 $-\bm{D}_d\dot{\bm{x}}$처럼 보인다.
이 힘을 실제 관절 effort로 바꾸기 위해 표~\ref{tab:jacobian-recap-for-impedance}의 자코비안 전치 관계를 사용한다.
\vmark{scope-s6-actual-joint-effort}%
여기서 $\bm{f}_{\mathrm{cmd}}$는 제어기가 로봇 끝단에 가하려는 힘이다.
로봇이 환경에 가하는 힘을 양수로 쓰는 표기에서는 작용-반작용 때문에 같은 접촉의 부호가 반대로 보일 수 있으므로, 센서 힘이나 벽 힘을 이 식에 넣기 전에 기준을 맞추어야 한다.
문손잡이를 미는 힘이 경첩 토크로 바뀌는 것처럼, 로봇 손끝에 만들고 싶은 힘도 각 관절에서 필요한 일반화 effort로 바뀌어야 한다.
\vmark{scope-s6-force-to-generalized}%
회전 관절만 놓고 보면 이 effort가 관절 토크이다.
절~\ref{sec:robotics-foundations}의 식~\eqref{eq:jacobian-transpose-foundation}이 그 유도이고, 이 절에서는 병진 자코비안 $\bm{J}_v$에 대해 그 결과만 사용한다.
따라서 이 단계에서는 자코비안을 거꾸로 푸는 것이 아니라, 손끝 힘이 각 관절에 요구하는 일반화 effort 기여분을 나누어 계산한다고 보면 된다.
관절이 $n$개이고 병진 힘만 다루면 $\bm{J}_v$는 보통 $3\times n$ 행렬이고, $\bm{J}_v^{\mathsf{T}}\bm{f}_{\mathrm{cmd}}$는 $n$개의 관절 effort를 만든다.
\vmark{scope-s6-n-joint-effort}%
이 절의 회전 관절 예제에서는 이 값들을 관절 토크로 읽는다.
2자유도 평면 팔에서는 이 말이 더 작게 보인다.
손끝 힘 \([f_x\ f_y]^{\mathsf{T}}\)가 두 관절 토크 \([\tau_1\ \tau_2]^{\mathsf{T}}\)로 바뀌며, 각 관절 토크는 해당 자코비안 열과 손끝 힘의 내적으로 정해진다.
\begin{equation}
  \bm{\tau}_{\mathrm{cmd}}
  =
  \bm{J}_v(\bm{q})^{\mathsf{T}}
  \bm{f}_{\mathrm{cmd}}.
  \label{eq:cartesian_to_joint_torque}
\end{equation}
이 식은 작업공간 힘을 만들기 위해 필요한 관절 토크 명령을 계산하는 출발점이지, 어떤 상황에서도 그 힘이 반드시 정확히 만들어진다는 보장은 아니다.
관절 토크가 포화되거나, 로봇이 특이점 근처에 있거나, 제어 주기가 느리거나, 모델 보상이 부족하면 실제 끝단 힘은 명령과 달라질 수 있다.
이 관계는 작업공간 제어와 운영공간 제어의 핵심 출발점이다 \cite{khatib1987operationalspace}.
운영공간 제어는 손끝 같은 작업공간 좌표에서 힘과 가속도를 직접 다루려는 로봇 제어 관점이다.
여기서는 그 전체 이론을 모두 쓰기보다, 작업공간 힘을 관절 토크로 옮기는 출발점만 사용한다.
만약 병진 힘뿐 아니라 회전 모멘트까지 포함한 6차원 렌치를 다루려면, 기하학적 자코비안 $\bm{J}_g$와 렌치 $\bm{w}$를 사용해 $\bm{\tau}=\bm{J}_g^{\mathsf{T}}\bm{w}$처럼 쓴다.
그림~\ref{fig:cartesian-impedance-block}은 지금까지의 직교 임피던스 흐름을 블록도로 요약한다.
작업공간에서는 위치와 속도 오차로 가상 힘을 만들고, 관절공간에서는 그 힘을 실제 모터가 낼 토크로 바꾸어 실행한다.
이 그림은 병진 위치 임피던스의 개념도이며, 전체 운영공간 역동역학을 모두 그린 것은 아니다.
실제 구현에서는 여기에 중력과 바이어스 보상, 토크 제한, 여유 자유도를 다루는 null-space 항이 더해질 수 있다.

\begin{figure}[!htbp]
  \centering
  \resizebox{\linewidth}{!}{%
  \begin{tikzpicture}[
      block/.style={draw=blue!55!black, rounded corners=2pt, fill=blue!4, align=center, minimum width=2.7cm, minimum height=1.06cm, font=\small, inner sep=5pt},
      calc/.style={draw=black!45, rounded corners=2pt, fill=black!3, align=center, minimum width=2.2cm, minimum height=0.86cm, font=\scriptsize, inner sep=4pt},
      input/.style={draw=black!45, rounded corners=2pt, fill=white, align=center, minimum width=1.85cm, minimum height=0.74cm, font=\scriptsize, inner sep=3pt},
      sum/.style={draw=black!65, circle, fill=white, minimum size=0.58cm, inner sep=0pt, font=\scriptsize},
      arr/.style={-{Latex[length=2.0mm]}, thick, draw=black!65},
      fb/.style={-{Latex[length=1.7mm]}, draw=black!45, thick},
      note/.style={draw=black!35, rounded corners=2pt, fill=black!3, align=center, font=\scriptsize, text width=5.9cm, inner sep=4pt}
    ]

    \node[input] (xd) at (0,1.05) {희망 위치\\$\bm{x}_d$};
    \node[input] (x) at (0,-1.05) {현재 위치\\$\bm{x}$};
    \node[sum] (sumx) at (1.65,0) {$\Sigma$};
    \node[calc] (err) at (3.25,0) {작업공간 오차\\$\bm{e}=\bm{x}-\bm{x}_d$};
    \node[input] (xdot) at (3.25,-1.45) {현재 속도\\$\dot{\bm{x}}$};
    \node[block, minimum width=3.2cm] (imp) at (6.15,0) {가상 스프링-댐퍼\\힘 생성\\$\bm{f}_{\mathrm{cmd}}=-\bm{K}_d\bm{e}-\bm{D}_d\dot{\bm{x}}$};
    \node[block, minimum width=2.7cm] (jt) at (9.55,0) {자코비안 전치\\토크 변환\\$\bm{\tau}_{\mathrm{cmd}}=\bm{J}_v(\bm{q})^{\mathsf T}\bm{f}_{\mathrm{cmd}}$};
    \node[block, minimum width=2.35cm] (robot) at (12.45,0) {로봇 팔\\환경과 접촉};
    \node[input] (fext) at (12.45,1.55) {외력\\$\bm{f}_{\mathrm{ext}}$\\환경 $\rightarrow$ 로봇};
    \node[calc] (state) at (12.45,-1.55) {센서/상태 추정\\$\bm{x},\dot{\bm{x}},\bm{q}$};

    \draw[arr] (xd) -| node[pos=0.28, above, font=\scriptsize] {$-$} (sumx);
    \draw[arr] (x) -| node[pos=0.28, below, font=\scriptsize] {$+$} (sumx);
    \draw[arr] (sumx) -- (err);
    \draw[arr] (err) -- (imp);
    \draw[arr] (xdot) -| (imp.south);
    \draw[arr] (imp) -- (jt);
    \draw[arr] (jt) -- (robot);
    \draw[arr] (fext) -- (robot);
    \draw[fb] (robot.south) -- (state.north);
    \draw[fb] (state.west) -- ++(-1.0,0) |- (x.east);
    \draw[fb] (state.west) -- ++(-1.0,0) |- (xdot.east);
    \draw[fb] (state.north west) -- ++(-0.85,0.48) |- (jt.south);

    \node[note] at (6.15,-2.55) {개념도: 병진 위치 임피던스 중심. 정지 목표 $\dot{\bm{x}}_d=\bm{0}$를 가정하고, 중력 보상과 토크 제한, 자세 임피던스는 생략했다.};

    \draw[densely dashed, blue!35] (1.2,-2.05) rectangle (7.95,1.95);
    \node[font=\scriptsize, text=blue!55!black] at (4.58,2.15) {작업공간에서 힘을 설계};
    \draw[densely dashed, green!35!black] (8.15,-2.05) rectangle (13.75,1.95);
    \node[font=\scriptsize, text=green!35!black] at (10.95,2.15) {관절 토크로 실행};
  \end{tikzpicture}%
  }
  \caption{직교 임피던스 제어의 기본 흐름. 작업공간에서 현재 위치와 희망 위치의 차이, 그리고 현재 속도로부터 가상 스프링-댐퍼 힘 $\bm{f}_{\mathrm{cmd}}$를 만든다. 이 병진 작업공간 힘은 가상일 또는 순간 일률 보존 관점에서 자코비안 전치 관계 $\bm{\tau}_{\mathrm{cmd}}=\bm{J}_v(\bm{q})^{\mathsf T}\bm{f}_{\mathrm{cmd}}$를 통해 관절 토크로 변환된다. 따라서 흐름은 운동 오차에서 힘 명령으로, 다시 관절 토크로 이어진다.}
  \label{fig:cartesian-impedance-block}
\end{figure}


직교 임피던스의 장점은 사람이 이해하기 쉽다는 것이다.
예를 들어 벽을 누르는 방향인 $x$축은 부드럽게 하고, 벽을 따라 미끄러지는 $y$축은 단단하게 만들 수 있다.
조립 작업에서는 삽입 방향은 부드럽게, 정렬 방향은 단단하게 만드는 식의 설계가 가능하다.
반면 단점도 있다.
자코비안, 특이점, 자세 변화, 관절 토크 제한을 고려해야 한다.
특히 로봇이 특이점 근처에 가면 작은 작업공간 속도나 위치 보정이 큰 관절 움직임을 요구할 수 있고, 역기구학과 운영공간 동역학 계산의 조건이 나빠진다.
여기서 조건이 나쁘다는 말은 작은 위치 오차, 센서 노이즈, 수치 오차가 계산 결과에서 크게 증폭될 수 있는 상태라는 뜻이다.
자코비안 전치식 $\bm{\tau}=\bm{J}_v^{\mathsf{T}}\bm{f}$ 자체는 작업공간 병진 힘을 관절 토크로 바꾸는 가상일 관계이지만, 특이점 근처에서는 원하는 작업공간 운동을 만들기 위한 역방향 계산과 토크 제한이 문제가 되기 쉽다.
$\bm{J}_v^{\mathsf{T}}\bm{f}$는 수식으로는 항상 계산할 수 있지만, 특이점 근처에서는 특정 힘이나 강성을 실제로 만들 수 있는 방향이 제한되고 필요한 관절 토크가 나빠질 수 있다.
따라서 힘 매핑 자체도 로봇 자세, 랭크, 토크 한계와 함께 읽어야 한다.
2링크 팔을 완전히 쭉 편 모습을 떠올리면 직관적이다.
팔이 일직선에 가까울수록 손끝을 어떤 방향으로는 조금 움직이기 위해 관절을 크게 움직여야 하거나, 아예 움직임 선택지가 줄어든다.
이런 자세 근처에서는 역기구학 계산이 예민해지고 조작성(manipulability) 지표가 작아진다.
조작성은 현재 자세에서 손끝을 여러 방향으로 얼마나 쉽게 움직일 수 있는지를 나타내는 지표이다.

\begin{table}[!htbp]
  \centering
  \caption{관절 임피던스와 직교 임피던스의 차이}
  \begin{tabular}{lll}
    \toprule
    구분 & 관절 임피던스 & 직교 임피던스 \\
    \midrule
    설계 좌표 & 관절각 $\bm{q}$ & 끝단 병진 위치 $\bm{x}$; 자세는 확장 모델 \\
    출력 직관 & 관절별 부드러움 & 공간 방향별 부드러움 \\
    장점 & 구현이 비교적 단순함 & 작업/접촉 방향 설계가 직관적 \\
    단점 & 끝단 강성이 자세에 따라 변함 & 자코비안, 특이점, 토크 제한 고려 필요 \\
    대표 용도 & 자세 유지, 안전한 관절 제어 & 접촉 작업, 표면 추종, 조립 \\
    \bottomrule
  \end{tabular}
\end{table}

\subsection{임피던스 제어, 어드미턴스 제어, 힘 제어}

처음 공부할 때는 세 제어법을 우열로 외우기보다, 로봇에게 실제로 명령할 수 있는 값이 무엇인지부터 보면 된다.
토크나 힘을 비교적 직접 만들 수 있는 로봇이라면 위치 오차에서 힘을 만드는 임피던스 제어가 자연스럽다.
반대로 산업용 로봇처럼 내부 위치 제어기가 강하고 외부에서는 목표 위치만 줄 수 있다면, 측정된 힘을 보고 목표 위치를 조금씩 움직이는 어드미턴스 제어가 현실적일 수 있다.
작업 자체가 표면을 일정한 힘으로 누르는 것이라면 힘 제어가 핵심이 되지만, 자유공간 이동과 접촉 안정성을 함께 보려면 위치/임피던스/힘 제어를 방향별로 조합하는 편이 많다.

임피던스 제어와 어드미턴스 제어는 서로 반대 방향의 관계를 설계한다.
물리적인 스프링-댐퍼 관계식 자체는 힘에서 운동으로도, 운동에서 힘으로도 읽을 수 있다.
다만 제어 구현 관점에서는 어느 값을 센서로 읽고 어느 값을 명령으로 내보내는지가 중요하다.
임피던스 제어는 위치와 속도 오차를 보고 힘 또는 토크 명령을 계산한다.
\begin{equation}
  \text{motion} \rightarrow \text{force}.
\end{equation}
예를 들어 위치 오차와 속도를 보고
\begin{equation}
  f_{\mathrm{cmd}}
  =
  -k_d e
  -
  d_d\dot{e}
\end{equation}
라는 힘 또는 토크를 만든다.
목표점이 정지해 있으면 $\dot{e}=\dot{x}$이므로 감쇠항이 속도에 비례한 브레이크처럼 보인다.
목표점에서 얼마나 벗어났는지, 얼마나 빠르게 움직이는지를 보고 그만큼 되돌리는 힘을 만드는 셈이다.

어드미턴스 제어(admittance control)는 반대이다.
외부 힘을 입력으로 받아 목표 운동을 만든다.
\begin{equation}
  \text{force} \rightarrow \text{motion}.
\end{equation}
예를 들어 힘 센서에서 측정한 외력 $f_{\mathrm{ext}}$를 가상의 질량-댐퍼-스프링 모델에 넣고, 기준 목표점에서 얼마나 물러날지 또는 따라갈지를 계산한다.
여기서 $f_{\mathrm{des}}$는 $f_{\mathrm{ext}}$와 같은 부호 규약으로 표현한 원하는 측정 외력이다.
예를 들어 환경이 로봇을 미는 힘을 양의 방향으로 읽기로 했다면, 원하는 힘도 같은 방향 기준으로 써야 한다.
반대로 로봇이 $+x$ 방향으로 벽을 미는데 벽이 로봇에 가하는 측정 힘이 $-x$ 방향이라면, 로봇이 벽을 $+10~\mathrm{N}$으로 밀기 바라는 상황은 이 부호 규약에서 $f_{\mathrm{des}}=-10~\mathrm{N}$으로 써야 한다.
기준 위치를 $x_d$라고 하고, 그 기준에서의 추가 변위를 $\Delta x$라고 쓰면
\begin{equation}
  m_d\Delta\ddot{x}
  +
  d_d\Delta\dot{x}
  +
  k_d\Delta x
  =
  f_{\mathrm{ext}}-f_{\mathrm{des}}
\end{equation}
와 같이 쓸 수 있다.
그다음 새로운 위치 명령을
\begin{equation}
  x_{\mathrm{cmd}}=x_d+\Delta x
\end{equation}
로 만든다.
이렇게 써야 외력이 0일 때 명령이 좌표 원점으로 끌려가는 것이 아니라, 원래 기준 목표점 $x_d$ 근처에 머문다는 의미가 분명해진다.
정적 직관만 보려면 더 간단하다.
목표 힘이 $f_{\mathrm{des}}=0$이고, 사람이 로봇을 $f_{\mathrm{ext}}=5~\mathrm{N}$로 민다고 하자.
가상의 강성을 $k_d=100~\mathrm{N/m}$로 두면 정상상태 추가 변위는
\begin{equation}
  \Delta x_{\mathrm{ss}}
  =
  \frac{f_{\mathrm{ext}}-f_{\mathrm{des}}}{k_d}
  =
  \frac{5}{100}
  =
  0.05~\mathrm{m}
\end{equation}
이다.
즉, 어드미턴스 제어기는 힘을 직접 없애려고 토크를 내는 대신, 위치 명령을 약 $5~\mathrm{cm}$ 움직여서 로봇이 밀리는 듯한 행동을 만들 수 있다.
이 부호가 맞는지 확인하는 간단한 테스트도 있다.
$f_{\mathrm{des}}=0$이고 양의 방향 외력이 로봇을 민다면, 위 식은 $\Delta x>0$을 만들어 명령 위치가 밀린 방향으로 움직인다.
만약 힘 센서가 환경이 로봇에 가한 힘이 아니라 로봇이 환경에 가한 힘을 양수로 보고한다면 측정값의 부호가 반대이므로, 측정 힘을 한 번 부호 반전하거나 힘 오차식을 반대로 써야 한다.
그리고 로봇의 내부 위치 제어기가 그 위치를 따라가게 한다.
토크 제어가 어렵고 위치 제어 인터페이스만 강하게 제공되는 산업용 로봇에서는 어드미턴스 제어가 유용한 경우가 많다.
힘 센서로 힘을 측정하고, 그 힘에 따라 목표 위치를 조금씩 움직여 주는 방식이기 때문이다.

힘 제어(force control)는 원하는 힘 자체를 추종하려는 제어이다.
예를 들어 표면을 $10~\mathrm{N}$으로 누르고 싶다면 힘 오차
\begin{equation}
  e_f=f_{\mathrm{des}}-f_{\mathrm{meas}}
\end{equation}
를 줄이도록 제어한다.
힘 제어는 연마, 누르기, 삽입, 표면 접촉처럼 접촉력이 직접 중요한 작업에서 필요하다.
하지만 순수 힘 제어는 자유공간에서는 목표가 애매해질 수 있고, 딱딱한 환경에서는 작은 위치 변화가 큰 힘 변화로 이어져 불안정해질 수 있다.
예를 들어 아무것도 없는 허공에서 $10~\mathrm{N}$으로 누르라고 하면, 무엇을 누르는지부터 정의되지 않는다.
접촉 대상이 없으면 힘 목표만으로는 어디로 가야 하는지 결정할 수 없다.
그래서 실제로는 위치 제어, 임피던스 제어, 힘 제어를 작업 방향에 따라 섞어 쓰는 경우가 많다.

하이브리드 위치/힘 제어(hybrid position/force control)는 이 섞어 쓰는 생각을 방향별로 정리한 접근이다.
예를 들어 표면을 닦는 작업에서는 표면에 수직인 방향으로는 원하는 누르는 힘을 맞추고, 표면을 따라가는 접선 방향으로는 위치나 속도 경로를 따라가게 할 수 있다.
구멍에 부품을 끼울 때도 삽입 방향과 옆 방향을 다르게 다루는 것이 자연스럽다.
수학적으로는 선택 행렬(selection matrix)로 어떤 축은 위치 제어, 어떤 축은 힘 제어에 맡긴다고 표현할 수 있지만, 처음에는 축별로 제어 목표를 나누는 방법이라고 생각하면 충분하다.
임피던스 제어는 이와 경쟁한다기보다, 각 방향에서 힘과 운동 사이의 관계를 얼마나 단단하거나 부드럽게 만들지를 정하는 또 다른 언어로 함께 쓰일 수 있다.

구현 관점에서 다시 정리하면, 임피던스 제어는 운동 오차를 힘으로 바꾸고, 어드미턴스 제어는 측정 힘을 운동 명령으로 바꾸며, 힘 제어는 측정 힘이 목표 힘에 가까워지도록 피드백을 닫는다.

\begin{table}[!htbp]
  \centering
  \caption{임피던스, 어드미턴스, 힘 제어의 개념 차이}
  \small
  \setlength{\tabcolsep}{4pt}
  \begin{tabularx}{\linewidth}{@{}>{\raggedright\arraybackslash}p{0.18\linewidth}>{\raggedright\arraybackslash}p{0.22\linewidth}>{\raggedright\arraybackslash}p{0.25\linewidth}>{\raggedright\arraybackslash}X@{}}
    \toprule
    방법 & 입력처럼 보는 것 & 출력처럼 만드는 것 & 좋은 상황 \\
    \midrule
    임피던스 제어 & 위치/속도 오차 & 힘/토크 & 토크 제어 가능, 접촉 반응 설계. 힘 목표 자체를 직접 추종하기보다 운동 오차에서 힘을 만든다 \\
    어드미턴스 제어 & 측정 외력 & 위치/속도 명령 & 위치 제어형 로봇, 힘 센서 사용 \\
    힘 제어 & 힘 오차 & 힘을 맞추는 동작 & 측정 힘 피드백을 닫아 $f_{\mathrm{meas}}$가 $f_{\mathrm{des}}$에 가까워져야 하는 작업 \\
    하이브리드 위치/힘 제어 & 방향별 위치 오차와 힘 오차 & 축마다 다른 명령 & 법선 방향은 힘, 접선 방향은 위치처럼 작업 방향별 목표가 다른 경우 \\
    \bottomrule
  \end{tabularx}
\end{table}

\begin{table}[!htbp]
  \centering
  \caption{상황별 제어 관점 선택의 출발점}
  \small
  \setlength{\tabcolsep}{4pt}
  \begin{tabularx}{\linewidth}{@{}>{\raggedright\arraybackslash}p{0.31\linewidth}>{\raggedright\arraybackslash}p{0.26\linewidth}>{\raggedright\arraybackslash}X@{}}
    \toprule
    상황 & 먼저 떠올릴 관점 & 이유 \\
    \midrule
    자유공간에서 정해진 경로를 반복 이동 & 위치 또는 궤적 제어 & 접촉이 없으면 힘-운동 관계보다 위치 정확도가 먼저 보인다 \\
    사람이 로봇을 밀면 부드럽게 물러나야 함 & 임피던스 또는 어드미턴스 제어 & 외력에 대한 변위와 속도 응답을 정해야 한다 \\
    토크 제어가 가능하고 접촉 반응을 직접 설계하고 싶음 & 임피던스 제어 & 위치/속도 오차를 원하는 힘 또는 토크로 바꾸기 쉽다 \\
    위치 명령만 받는 로봇에 힘 센서가 달려 있음 & 어드미턴스 제어 & 측정 힘을 목표 위치 변화로 바꾸어 내부 위치 제어기에 넘길 수 있다 \\
    일정한 누르는 힘이 작업 품질을 좌우함 & 힘 제어 또는 하이브리드 제어 & 접촉 방향의 힘 오차를 직접 줄여야 한다 \\
    \bottomrule
  \end{tabularx}
\end{table}

\subsection{목적에 따라 어떤 파라미터를 조절해야 하는가?}

임피던스 튜닝은 큰 값을 찾는 일이 아니다.
먼저 안전한 힘과 속도 한계를 정하고, 작은 값에서 시작해 플롯을 보며 하나씩 올리는 과정에 가깝다.
같은 강성 증가도 자세, 관절 토크, 속도 제한, 내부 위치 제어 대역폭, 센서 노이즈, 자코비안 조건수에 따라 전혀 다른 결과로 나타날 수 있기 때문이다.
특히 특이점이나 관절 제한 근처에서는 같은 작업공간 강성도 훨씬 큰 관절 속도나 토크 요구로 바뀔 수 있다.
이 절은 실험자가 플롯을 보며 어떤 질문을 던지고, 그 질문에 맞춰 어떤 파라미터를 움직이는지 정리한다.

외력이 있을 때 최종 위치 오차를 줄이고 싶다면 강성 $k_d$를 키운다.
여기서 $F_0$는 환경이 로봇에 가하는 힘을 기준으로 쓴 값이고, 로봇이 환경에 가하는 힘은 반대 방향이다.
정상상태에서 부호를 포함하면 $e_{\mathrm{ss}}=F_0/k_d$이고, 크기로 보면 $|e_{\mathrm{ss}}|=|F_0|/k_d$이므로, 같은 외력 크기에 대해 위치 오차의 크기가 줄어든다.
하지만 강성을 너무 크게 하면 접촉 순간 힘이 급격히 커지고, 센서 노이즈나 모델 오차에도 민감해진다.
토크 포화가 생기면 오히려 응답이 나빠질 수 있다.

도달 속도를 빠르게 하고 싶다면, 먼저 어떤 질량을 기준으로 말하는지 구분해야 한다.
1축 목표 임피던스 모델이 실제 닫힌루프에서 잘 구현된다고 가정하면
\begin{equation}
  \omega_n=\sqrt{\frac{k_d}{m_d}}
\end{equation}
이다.
여기서 $m_d$는 제어기가 목표로 삼는 희망 관성이다.
토크 제어와 동역학 보상이 충분해서 이 희망 관성이 잘 구현될 때는, $m_d$가 고정된 상태에서 $k_d$를 키우면 목표 모델의 고유진동수가 커지고 응답이 빨라진다.
다만 실제 로봇에서 관찰되는 속도와 진동 주기는 현재 자세와 방향의 끝단 유효질량, 내부 위치 제어 대역폭, 포화, 지연의 영향을 함께 받는다.
교육용 데모에서 $m_d=1~\mathrm{kg}$ 또는 $m_{\mathrm{nom}}=1~\mathrm{kg}$ 같은 명목 질량으로 $d_d$를 계산했다면, 앞에서 말한 대로 그 감쇠비는 물리 감쇠비가 아니라 튜닝 기준이다.
따라서 먼저 $k_d$를 조금씩 올리고, 선택한 감쇠비가 유지되도록 내부 감쇠계수 $d_d$를 다시 계산한 뒤, 실제 플롯에서 속도, 오버슈트, actuator effort, 즉 구동기 부담 지표를 확인한다.
하지만 빠른 응답은 충분한 감쇠가 함께 있어야 한다.
강성만 키우고 감쇠를 그대로 두면 감쇠비가 낮아져 진동과 오버슈트가 커질 수 있다.

진동을 줄이고 싶다면 먼저 1축 선형 모델의 감쇠비 $\zeta$를 출발점으로 조절한다.
$\zeta \approx 0.7$은 빠른 과도응답의 대표값이고, $\zeta \approx 1$은 오버슈트를 줄이는 대표값이다.
실제 제어기 내부에서는 식 \eqref{eq:damping_from_ratio}처럼 선택한 $\zeta$와 $k_d$, $m_d$로부터 감쇠계수 $d_d$를 계산하는 식으로 구현할 수 있다.
다만 접촉 상황에서 더 큰 $\zeta$나 $d_d$가 항상 더 안전한 것은 아니다.
감쇠계수 $d_d$가 너무 커지면 움직임이 끈적하고 둔해지며, 속도 측정 노이즈, 지연, 포화, 구동기 effort 피크가 커질 수 있다.

적절한 접촉력을 만들고 싶다면 강성과 목표점 위치를 함께 봐야 한다.
벽이 있고 로봇의 가상 평형점을 벽 안쪽으로 목표 오프셋 $\delta_d$만큼 일부러 넣는다고 하자.
이 $\delta_d$는 로봇이 실제로 벽을 뚫고 들어간 거리가 아니라, 제어기가 목표점을 벽 뒤에 둔 설계량이다.
강체에 가까운 벽, 거의 0인 속도, 포화되지 않은 제어 입력, 충분한 제어 대역폭을 가정하고, 손끝이 벽 표면 근처에 머물러 실제 오차 크기가 목표 오프셋과 비슷하다면 접촉력은 대략
\begin{equation}
  |f_{\mathrm{contact}}|\approx k_d\delta_d
\end{equation}
이다.
예를 들어 목표점을 벽 안쪽으로 $\delta_d=2~\mathrm{cm}$ 넣고 $k_d=500~\mathrm{N/m}$로 잡으면, 정지 근처에서는 약 $10~\mathrm{N}$의 접촉력을 기대할 수 있다.
따라서 같은 접촉력을 만들려면 큰 강성에서는 작은 목표 오프셋이 필요하고, 작은 강성에서는 큰 목표 오프셋이 필요하다.
큰 강성은 위치 오차가 작고 단단하지만 힘 변화가 날카롭다.
작은 강성은 부드럽고 접촉력 변화가 완만할 수 있지만 목표 위치에서 더 많이 밀린다.
안전성은 강성 하나가 아니라 속도 제한, 힘 제한, 감쇠, 환경 강성, 제어 대역폭까지 함께 좌우한다.
만약 환경도 유한한 강성을 가진다면 로봇의 가상 스프링과 환경 스프링이 함께 변형되므로, 접촉력은 단순히 $k_d\delta_d$ 하나로만 정해지지 않는다.
목표 오프셋 전체가 로봇 쪽 스프링 변형으로만 들어가는 것이 아니라, 일부는 로봇의 위치 오차로, 일부는 환경의 변형으로 나뉜다.
그래서 부드러운 물체를 누를 때는 같은 목표 오프셋을 주어도 벽이 강체일 때보다 힘이 작게 나오거나, 변형 분포가 전혀 다르게 보일 수 있다.
이때는 로봇 강성, 환경 강성, 목표 오프셋이 직렬로 연결된 두 스프링 문제처럼 함께 힘의 균형을 이룬다고 생각하면 좋다.
이 균형을 수식 한 줄로 적어 보자.
환경을 강성 $k_{\mathrm{env}}$의 스프링으로 보면, 정지 근처 평형에서는 하나의 접촉력 $f$가 로봇 가상 스프링과 환경 스프링을 동시에 누르므로
\[
  f
  =
  k_d\,\delta_{\mathrm{robot}}
  =
  k_{\mathrm{env}}\,\delta_{\mathrm{env}}
\vmark{contact-series-stiffness-split}%
\]
이고, 목표 오프셋은 두 변형으로 나뉘어 $\delta_{\mathrm{robot}}+\delta_{\mathrm{env}}\approx\delta_d$가 된다.
첫 식에서 $\delta_{\mathrm{robot}}=f/k_d$, $\delta_{\mathrm{env}}=f/k_{\mathrm{env}}$로 풀어 합 식에 넣으면 $f(1/k_d+1/k_{\mathrm{env}})\approx\delta_d$이므로
\[
  f
  \approx
  \frac{k_d\,k_{\mathrm{env}}}{k_d+k_{\mathrm{env}}}\,\delta_d.
\]
즉 직렬 합성 강성은 두 강성 중 작은 쪽에 끌려간다.
숫자로 보면, 위의 $k_d=500~\mathrm{N/m}$, $\delta_d=2~\mathrm{cm}$ 예제에서 벽이 강체에 가까우면($k_{\mathrm{env}}\to\infty$) $f\to k_d\delta_d=10~\mathrm{N}$이지만, 부드러운 물체 $k_{\mathrm{env}}=500~\mathrm{N/m}$이면 합성 강성이 $250~\mathrm{N/m}$라 같은 목표 오프셋에서 $f\approx5~\mathrm{N}$으로 절반이 된다.
\vmark{contact-series-stiffness-numeric}%

접촉 튜닝에서는 기호도 한 번 정리해 두는 편이 좋다.
\begin{table}[!htbp]
  \centering
  \caption{접촉 예제에서 헷갈리기 쉬운 강성, 오프셋, 힘 기호}
  \label{tab:contact-symbols}
  \small
  \setlength{\tabcolsep}{4pt}
  \begin{tabularx}{\linewidth}{@{}>{\raggedright\arraybackslash}p{0.18\linewidth}>{\raggedright\arraybackslash}p{0.28\linewidth}>{\raggedright\arraybackslash}X@{}}
    \toprule
    기호 & 의미 & 처음 읽을 때의 해석 \\
    \midrule
    $k_d$ & 로봇 쪽 목표 임피던스의 가상 스프링 강성 & 손끝이 목표점에서 밀렸을 때 얼마나 단단하게 되돌아가려 하는지를 정한다 \\
    $\delta_d$ & 벽 안쪽으로 일부러 넣은 목표점 오프셋 & 실제 침투가 아니라, 정적 접촉력을 만들기 위한 설계 여유이다 \\
    $k_{\mathrm{wall}}$ & Lab04 결정론적 가상 벽 모델의 벽 강성 & 로봇 제어 강성이 아니라, 교육용 벽 신호를 계산하는 별도 파라미터이다 \\
    $\delta$ & $\max(0,x-x_{\mathrm{wall}})$로 계산한 침투 깊이 & 손끝이 가상 벽 경계 안쪽으로 얼마나 들어갔는지를 나타낸다 \\
    \texttt{force\_virtual} & $k_{\mathrm{wall}}$, $\delta$, 접근 속도 감쇠항으로 만든 계산 신호 & 힘 센서가 잰 실제 접촉력이 아니라, 플롯을 읽기 위한 교육용 벽 힘이다 \\
    \bottomrule
  \end{tabularx}
\end{table}

따라서 MuJoCo의 실제 접촉 해석기가 계산하는 제약력과, 뒤에서 설명하는 결정론적 가상 벽의 벽 힘은 같은 물리량으로 취급하면 안 된다.
가상 벽 데모에서는 정해 둔 벽 강성 $k_{\mathrm{wall}}$과 실제 침투 깊이 $\delta$로부터 $|f_{\mathrm{wall}}|\approx k_{\mathrm{wall}}\delta$라는 별도의 교육용 힘을 만든다.
목표 후퇴와 DLS 오프셋은 실제 침투 깊이를 다시 바꿀 수 있으므로, 이 값을 $k_d\delta_d$와 같은 물리량처럼 읽으면 안 된다.

마지막으로, 빠름과 부드러움을 완전히 따로 고를 수는 없다는 점도 기억해야 한다.
1축 목표 임피던스 모델에서는
\begin{equation}
  k_d=m_d\omega_n^2
\end{equation}
이므로 $k_d$, $m_d$, $\omega_n$을 모두 따로따로 마음대로 고르는 것이 아니다.
셋 중 둘을 정하면 나머지 하나는 정해진다.
교육용 데모처럼 $m_{\mathrm{nom}}$을 고정해 두면 $k_d$를 낮추는 순간 더 부드러워질 뿐 아니라 $\omega_n$도 낮아져 느려지는 경향이 있다.
반대로 같은 정적 부드러움 $k_d$를 유지하면서 더 빠른 $\omega_n$을 원하면 더 작은 희망 관성 $m_d$를 가정해야 하는데, 실제 로봇이나 position actuator 구현에서는 그 관성을 마음대로 만들 수 없을 수 있다.
또 $\zeta$를 정하면 식 \eqref{eq:damping_from_ratio}을 통해 $d_d$도 함께 정해지므로, $d_d$와 $\zeta$ 역시 서로 독립된 두 손잡이가 아니다.
그래서 처음 튜닝할 때의 순서는 먼저 정적 부드러움으로 $k_d$를 잡고, 그다음 $\zeta$로 흔들림을 줄인 뒤, 실제 플롯에서 속도와 effort가 허용 범위 안에 있는지 확인하는 쪽이 자연스럽다.

처음 튜닝할 때는 다음 순서가 가장 안정적인 출발점이다.
강성 숫자부터 외우기보다, 먼저 접촉 방향에서 어느 정도 힘을 받을 때 얼마나 밀려도 되는지 정한다.
예를 들어 접촉 방향에서 $10~\mathrm{N}$을 받을 때 $2~\mathrm{cm}$ 정도 밀려도 된다고 정하면, 정적 근사에서는
\begin{equation}
  k_d \approx \frac{10~\mathrm{N}}{0.02~\mathrm{m}}
  =
  500~\mathrm{N/m}
\end{equation}
가 출발값이 된다.
같은 $10~\mathrm{N}$을 더 부드럽게 만들고 싶다면 $k_d$만 낮추는 것이 아니라, 허용 변위나 목표 오프셋도 함께 다시 계산해야 한다.
그다음 실행해 보고, 최종 오차는 괜찮은데 도착 과정에서 많이 흔들리면 감쇠비를 올린다.
반대로 움직임이 너무 둔하고 effort가 커지면 감쇠비나 강성을 다시 낮춘다.
최종 위치로 가는 동안 흔들림을 얼마나 허용할지 정하기 위해 $\zeta$를 고르고, 내부에서는 $d_d=2\zeta\sqrt{m_dk_d}$ 또는 명목 질량을 쓰는 근사식으로 감쇠계수를 계산한다.
예를 들어 $m_{\mathrm{nom}}=1~\mathrm{kg}$, $\zeta=0.7$, $k_d=500~\mathrm{N/m}$를 쓰면
\begin{equation}
  d_d
  =
  2\zeta\sqrt{m_{\mathrm{nom}}k_d}
  =
  2\cdot0.7\sqrt{1\cdot500}
  \approx
  31.3~\mathrm{N\,s/m}
\end{equation}
이다.
이 숫자 예제에서 기대하는 정적 위치 오차는 $10~\mathrm{N}/500~\mathrm{N/m}=0.02~\mathrm{m}$이고, 부족감쇠 기준 진동 빠르기는 $\omega_d\approx \omega_n\sqrt{1-0.7^2}$처럼 $\omega_n$보다 조금 낮다.
플롯에서는 최종 위치 오차가 대략 $2~\mathrm{cm}$ 근처로 가는지, 그 과정에서 오버슈트가 얼마나 남는지, 그리고 effort 지표가 포화에 닿는지 함께 확인한다.
Lab04에서는 이 $d_d$가 실제 끝단 물리 감쇠계수라는 보증값이 아니라, position actuator와 DLS 목표 오프셋 위에서 응답을 조절하기 위한 내부 기준값임을 다시 확인해야 한다.
이 지점을 놓치면 이상적인 임피던스 제어식과 Lab04의 교육용 구현을 같은 것으로 읽기 쉽다.
표~\ref{tab:impedance-implementation-ladder}처럼 구현 층위를 나누어 보면 구분이 선명해진다.

\begin{table}[!htbp]
  \centering
  \caption{이상적인 임피던스 식에서 Lab04 관찰까지의 구현 사다리}
  \label{tab:impedance-implementation-ladder}
  \small
  \setlength{\tabcolsep}{4pt}
  \begin{tabularx}{\linewidth}{@{}>{\raggedright\arraybackslash}p{0.22\linewidth}>{\raggedright\arraybackslash}p{0.27\linewidth}>{\raggedright\arraybackslash}p{0.22\linewidth}>{\raggedright\arraybackslash}X@{}}
    \toprule
    층위 & 계산하는 것 & 실제 명령 & 처음 읽을 때의 해석 \\
    \midrule
    목표 임피던스 모델 & $\bm{M}_d$, $\bm{D}_d$, $\bm{K}_d$가 만드는 힘-운동 관계 & 아직 특정 하드웨어 명령은 아니다 & 만들고 싶은 접촉 행동의 목표식이다 \\
    토크 제어 구현 & $\bm{f}_{\mathrm{cmd}}$와 $\bm{\tau}=\bm{J}_v^{\mathsf{T}}\bm{f}_{\mathrm{cmd}}$ & 관절 토크 또는 힘/토크 명령 & 동역학 보상과 구동기 대역폭이 충분할 때 이상식에 가까워진다 \\
    위치 명령형 근사 & 목표 위치 보정, 목표 궤적 필터, 또는 목표 오프셋 & 관절 또는 끝단 위치 명령 & 실제 힘을 직접 명령하지 않고 임피던스적인 거동을 흉내 낸다 \\
    Lab04 가상 벽 데모 & \texttt{force\_virtual}, 목표 후퇴량, DLS 관절 목표 오프셋 & \texttt{ctrl} 형태의 관절 목표 위치 & 실제 접촉력 검증이 아니라 계산된 벽 신호와 effort 지표로 강성/감쇠 효과를 관찰한다 \\
    \bottomrule
  \end{tabularx}
\end{table}

따라서 Lab04는 $\bm{\tau}=\bm{J}_v^{\mathsf{T}}\bm{f}$를 그대로 실현하는 완전한 운영공간 임피던스 검증이 아니다.
\vmark{scope-s6-lab04-caveat}%
대신 같은 강성-감쇠 직관을 position actuator, DLS 목표 오프셋, 계산된 가상 벽 신호 위에서 관찰하게 만든 교육용 구현이다.
즉, $k_d$는 힘을 받을 때 얼마나 밀려도 되는지를, $\delta_d$는 정적 접촉에서 얼마나 누르고 싶은지를, $\zeta$는 가는 동안 얼마나 출렁여도 되는지를 정하는 값이다.
토크나 힘을 직접 명령할 수 있으면 이 관계를 임피던스 제어의 힘/토크 명령으로 구현한다.
위치 명령만 줄 수 있으면 측정 또는 추정 외력으로 목표 위치를 바꾸는 어드미턴스식 외부 루프를 둘 수 있다.
엄밀한 어드미턴스 제어는 측정 또는 추정 외력을 입력으로 받아 운동 명령을 만드는 motion generator이다.
Lab04 가상 벽은 실측 힘 대신 $\delta$와 $\dot{\delta}_+$에서 계산한 결정론적 벽 신호로 목표점을 후퇴시키는 교육용 대리 구현에 해당한다.

\begin{table}[!htbp]
  \centering
  \caption{처음 튜닝할 때의 임피던스 파라미터 선택 순서}
  \small
  \setlength{\tabcolsep}{4pt}
  \begin{tabularx}{\linewidth}{@{}>{\raggedright\arraybackslash}p{0.25\linewidth}>{\raggedright\arraybackslash}p{0.35\linewidth}>{\raggedright\arraybackslash}X@{}}
    \toprule
    먼저 묻는 질문 & 예시 & 출발 설정 \\
    \midrule
    어느 방향을 조절할까? & 벽을 누르는 $x$ 방향 & 해당 방향의 $k_d$와 $\zeta$부터 조절한다. \\
    얼마나 밀려도 될까? & $10~\mathrm{N}$에서 $2~\mathrm{cm}$ 허용 & $k_d \approx F/\Delta x = 500~\mathrm{N/m}$를 출발값으로 둔다. \\
    같은 힘을 더 부드럽게 만들까? & $10~\mathrm{N}$에서 $5~\mathrm{cm}$ 허용 & $k_d$를 낮추고 $\delta_d \approx F/k_d$를 다시 계산한다. \\
    흔들림은 어느 정도 허용할까? & 빠르지만 약간의 오버슈트 허용, 또는 오버슈트 최소화 & $\zeta\approx0.7$ 또는 $\zeta\approx1.0$에서 시작하고 $d_d$를 다시 계산한다. \\
    실제 로봇에 무엇을 명령할 수 있나? & 토크 명령 가능, 또는 위치 명령만 가능 & 임피던스 구현인지 어드미턴스 구현인지 선택한다. \\
    \bottomrule
  \end{tabularx}
  \label{tab:beginner-impedance-tuning-order}
\end{table}

작업 방향마다 다른 성질을 주고 싶다면 강성 행렬을 방향별로 다르게 설정한다.
예를 들어 표면을 따라 움직이는 작업에서는 표면 법선 방향은 낮은 강성으로 부드럽게 하고, 표면 접선 방향은 높은 강성으로 경로를 잘 따라가게 할 수 있다.
구멍에 부품을 끼우는 작업에서는 삽입 방향은 비교적 부드럽게, 옆으로 틀어지는 방향은 더 단단하게 만들 수 있다.
이런 설계는 직교 임피던스 제어에서 특히 직관적이다.

\begin{table}[!htbp]
  \centering
  \caption{목적별 파라미터 조절 직관}
  \small
  \setlength{\tabcolsep}{4pt}
  \begin{tabularx}{\linewidth}{@{}>{\raggedright\arraybackslash}p{0.22\linewidth}>{\raggedright\arraybackslash}p{0.25\linewidth}>{\raggedright\arraybackslash}X@{}}
    \toprule
    목적 & 주로 조절할 값 & 주의점 \\
    \midrule
    위치 오차 감소 & 강성 $k_d$ 증가 & 접촉 힘, 토크 포화, 노이즈 증가 \\
    빠른 도달 & $\omega_n$ 증가, 보통 $k_d$ 증가 & 감쇠도 함께 증가 필요 \\
    진동 감소 & 감쇠비 $\zeta$ 증가 & 너무 크면 둔하고 끈적함 \\
    부드러운 접촉 & 접촉 방향 $k_d$ 감소 & 위치 오차와 변위 증가 \\
    큰 정적 접촉력 & $k_d$ 또는 $\delta_d$ 증가 & 강체 접촉 근사에서만 단순히 읽고, Lab04 가상 벽 힘과 구분한다 \\
    표면 추종 & 방향별 $\bm{K}_d$, $\bm{D}_d$ 설정 & 좌표계와 자코비안 확인 필요 \\
    \bottomrule
  \end{tabularx}
\end{table}

실험을 하다 보면 보통 목표가 아니라 현상이 먼저 보인다.
그래서 실제 튜닝은 다음처럼 관찰에서 출발하는 편이 더 쉽다.

\begin{table}[H]
  \centering
  \caption{관찰된 현상에서 시작하는 임피던스 튜닝}
  \small
  \setlength{\tabcolsep}{4pt}
  \begin{tabularx}{\linewidth}{@{}>{\raggedright\arraybackslash}p{0.25\linewidth}>{\raggedright\arraybackslash}p{0.29\linewidth}>{\raggedright\arraybackslash}p{0.21\linewidth}>{\raggedright\arraybackslash}X@{}}
    \toprule
    관찰된 현상 & 먼저 의심할 것 & 조절 방향 & 함께 볼 신호 \\
    \midrule
    최종 위치 오차가 크다 & 강성이 낮거나 토크가 이미 포화됨 & $k_d$를 조금 올리되 포화 확인 & 위치 오차, 외력, 구동기 effort \\
    도달이 느리다 & 고유진동수 또는 강성이 낮음 & $k_d$를 올리고 $\zeta$ 유지 & 상승시간, 속도, 토크 제한 \\
    오버슈트나 진동이 크다 & 감쇠비가 낮거나 강성이 너무 큼 & $\zeta$를 올리고 필요하면 $k_d$ 완화 & 속도, 오버슈트, 진동 주기 \\
    움직임이 둔하고 끈적하다 & 감쇠계수 $d_d$가 과함 & $\zeta$를 낮추거나 목표 속도를 완만하게 조정 & 정착시간, 속도 노이즈 \\
    접촉력이 너무 크다 & 접촉 방향 강성 또는 목표 오프셋이 큼 & $k_d$ 또는 $\delta_d$를 낮춤 & 접촉력, 침투 깊이, 토크 포화 \\
    같은 힘을 더 부드럽게 만들고 싶다 & 큰 강성과 작은 변위 조합 & $k_d$를 낮추고 $\delta_d$를 다시 계산 & 힘 변화율, 저장 에너지, 변위 \\
    특정 방향만 불안정하다 & 자코비안 조건수, 방향별 강성, 관절 제한 & 해당 방향 $\bm{K}_d$, $\bm{D}_d$를 낮춤 & 조건수, joint speed, 구동기 effort \\
    \bottomrule
  \end{tabularx}
\end{table}

실험을 바꿀 때는 이번 실행의 질문을 하나만 정한다.
여기서 하나씩 바꾼다는 말은 숫자 하나만 절대적으로 바꾸라는 뜻은 아니다.
예를 들어 강성을 올리면서 같은 감쇠비를 유지하려면 $d_d$를 다시 계산할 수 있고, 같은 정적 접촉력 목표를 유지하려면 $\delta_d$도 함께 조정할 수 있다.
중요한 것은 이번 실행에서 검증할 가설을 하나로 유지하는 것이다.
더 단단하게 만들 것인지, 더 감쇠시킬 것인지, 포화를 줄일 것인지가 한 번에 섞이면 결과가 좋아져도 왜 좋아졌는지 알기 어렵다.
먼저 안전한 힘과 토크 한계를 정하고, 목표점을 갑자기 점프시키기보다 부드럽게 이동시키며, 위치 오차, 속도, 접촉력, 구동기 effort, 저장 에너지를 함께 보아야 한다.
뒤의 MuJoCo 랩은 바로 이 표를 실제 관찰 절차로 바꾼 것이다.
설정 파일에서 하나를 바꾸고, 위치, 속도, 힘, 제어 입력 플롯에서 어떤 응답 지표가 움직였는지 확인한다.

이제 두 그림으로 이 직관을 고정해 두자.
가상 퍼텐셜 에너지는 강성을 목표점 주변의 에너지 지형으로 읽게 해주고, 가상 벽은 그 지형이 접촉 상황에서 힘과 침투 깊이로 어떻게 보이는지 보여준다.
다음 두 절을 읽을 때는 새 공식을 외우기보다, 강성을 바꾸면 에너지 지형의 기울기가 어떻게 달라지고 감쇠를 바꾸면 벽에 닿는 순간의 튕김이 어떻게 달라지는지를 보면 된다.

\subsection{가상 퍼텐셜 에너지와 감쇠}

앞에서 같은 $10~\mathrm{N}$이라도 강성과 변위가 다르면 다른 물리 상황이라고 했다.
이번에는 그 차이를 목표점 주변의 에너지 지형으로 다시 해석한다.
식 \eqref{eq:virtual_spring_damper}의 강성항은 가상의 퍼텐셜 에너지에서 유도할 수 있다.
먼저 1차원부터 보자.
행렬 용어가 낯설다면 이 절에서는 1차원 식과 표만 먼저 읽어도 충분하다.
헤시안, 양의 정부호, 고유값 같은 말은 방향별 강성을 더 깊게 읽고 싶을 때 필요한 수학적 이름이다.
헤시안은 에너지 그릇이 얼마나 휘어 있는지 보는 곡률 행렬이고, 양의 정부호는 모든 방향으로 조금 움직이면 에너지가 증가한다는 조건이며, 고유값은 그 방향별 가파름을 나타내는 숫자라고 생각하면 된다.
에너지 함수
\begin{equation}
  U(e)=\frac{1}{2}k_de^2
\end{equation}
이 에너지는 목표점 $e=0$에서 가장 낮고, 양쪽으로 멀어질수록 커진다.
목표점이 최저점인 위로 열린 에너지 그릇에서 낮은 곳으로 내려가듯, 스프링 힘은 에너지가 줄어드는 목표점 쪽으로 향한다.
즉 1차원에서는 $k_d>0$일 때 목표점이 고립된 최소점이 된다.
행렬형에서는 에너지 함수의 곡률을 나타내는 헤시안이 $\nabla^2U=\bm{K}_d$이다.
따라서 모든 방향에서 목표점으로 돌아오게 만들려면 $\bm{K}_d$가 대칭 양의 정부호여야 한다.
행렬 표현은 이 1차원 그릇을 $x$, $y$, $z$ 방향으로 확장한 표현이다.
위치 오차를
\begin{equation}
  \bm{e}=\bm{x}-\bm{x}_d
\end{equation}
라고 하면, 목표점 주변의 가상 스프링 에너지는
\begin{equation}
  U(\bm{x})
  =
  \frac{1}{2}\bm{e}^{\mathsf{T}}\bm{K}_d\bm{e}
\end{equation}
로 둘 수 있다.
이때 에너지 함수로 해석하려면 강성 행렬은 대칭이어야 한다.
모든 방향을 닫힌 그릇처럼 만들려면 보통 $\bm{K}_d=\bm{K}_d^{\mathsf{T}}\succ0$를 가정한다.
반면 어떤 방향을 의도적으로 자유롭게 열어 두는 교육용 예제나 표면 추종 문제에서는 $\bm{K}_d\succeq0$인 준정부호 강성을 쓸 수 있다.
이 경우 $U\ge0$이지만, 영강성 방향에서는 최소점이 평평한 골짜리로 이어져 그 방향의 복원력은 없다.

\begin{table}[!htbp]
  \centering
  \caption{강성 크기를 에너지 지형으로 읽는 법}
  \small
  \begin{tabularx}{\linewidth}{@{}>{\raggedright\arraybackslash}p{0.23\linewidth}>{\raggedright\arraybackslash}p{0.33\linewidth}>{\raggedright\arraybackslash}X@{}}
    \toprule
    강성 & 에너지 지형 & 힘의 감각 \\
    \midrule
    0 & 그 방향으로 바닥이 평평하게 열린다 & 돌아오려는 힘이 없다 \\
    작음 & 완만한 그릇 & 많이 밀리지만 부드럽게 돌아온다 \\
    큼 & 가파른 그릇 & 조금만 벗어나도 큰 힘으로 돌아온다 \\
    \bottomrule
  \end{tabularx}
\end{table}

$\succeq0$이면 모든 방향에서 에너지가 음수가 되지 않는다.
쉽게 말해 어느 방향으로 밀어도 음의 에너지가 생기는 이상한 지형이 아니라, 평평하거나 위로 올라가는 그릇 모양을 만든다는 뜻이다.
$\succ0$이면 목표점이 유일한 최소점이고, 어떤 방향의 강성이 0이면 그 방향은 닫힌 그릇이 아니라 열린 골짜기처럼 남는다.
따라서 매우 작은 강성과 정확히 0인 강성은 구분해야 한다.
매우 작은 강성은 약하게나마 돌아오려는 힘을 만들지만, 0 강성은 그 방향으로는 돌아오려는 힘을 만들지 않는다.
일반적으로 방향별 강성은 $\bm{K}_d$의 고유벡터와 고유값으로 해석한다.
고유벡터는 주된 강성 방향을, 고유값은 그 방향의 강성 크기를 뜻한다.
대각 행렬은 좌표축이 이 주된 강성 방향과 일치하는 특별히 쉬운 경우이다.
대각 행렬로 쓰면 더 직관적이다.
\begin{equation}
  U(\bm{x})
  =
  \frac{1}{2}k_x e_x^2
  +
  \frac{1}{2}k_y e_y^2
  +
  \frac{1}{2}k_z e_z^2.
\end{equation}
즉, $x$, $y$, $z$ 방향에 서로 다른 가상의 스프링을 달아 놓은 것과 같다.
힘은 퍼텐셜 에너지의 음의 기울기이므로
\begin{equation}
  \bm{f}_k
  =
  -\frac{\partial U}{\partial \bm{x}}
  =
  -\bm{K}_d(\bm{x}-\bm{x}_d)
\end{equation}
가 된다.
즉, 강성 행렬 $\bm{K}_d$는 목표점 주변에 어떤 모양의 에너지 지형을 만들 것인지를 정한다.
큰 강성은 가파른 에너지 지형을 만들고, 작은 위치 오차에도 큰 복원력을 발생시킨다.
그림~\ref{fig:virtual-potential-stiffness}은 이 관계를 오차 좌표에서 그린 것이다.
등고선은 같은 힘이 아니라 같은 가상 퍼텐셜 에너지 $U$를 뜻한다.
힘은 이 에너지 지형의 음의 기울기 방향으로 생긴다.
등방 강성에서는 기울기가 목표점 방향과 나란하지만, 방향별 강성이 다르면 힘 화살표가 항상 목표점을 정확히 정조준하지는 않는다.
그래서 방향별 강성은 힘의 크기뿐 아니라 에너지 지형과 복원력 방향도 바꾼다.
숫자로 보면 더 분명하다.
위치 오차가
\begin{equation}
  \bm{e}
  =
  \begin{bmatrix}
  0.01 & 0.01
  \end{bmatrix}^{\mathsf{T}}~\mathrm{m}
\end{equation}
이고,
\begin{equation}
  \bm{K}_d
  =
  \begin{bmatrix}
  100 & 0 \\
  0 & 1000
  \end{bmatrix}
  ~\mathrm{N/m}
\end{equation}
이라면 스프링 힘은
\begin{equation}
  \bm{f}_k
  =
  -\bm{K}_d\bm{e}
  =
  \begin{bmatrix}
  -1 & -10
  \end{bmatrix}^{\mathsf{T}}~\mathrm{N}
\end{equation}
이다.
오차는 $x$와 $y$ 방향으로 똑같이 $1~\mathrm{cm}$였지만, $y$ 방향 강성이 10배 크기 때문에 힘은 대부분 $y$ 방향으로 생긴다.
따라서 방향별 강성을 다르게 주면 힘이 단순히 목표점을 향한 직선 방향으로만 생기지 않고, 더 단단한 방향의 기울기가 크게 반영된다.

\begin{figure}[!htbp]
  \centering
  \resizebox{0.96\linewidth}{!}{%
  \begin{tikzpicture}[
      axis/.style={-{Latex[length=2.0mm]}, draw=black!55},
      contour/.style={draw=blue!45, line width=0.55pt},
      softcontour/.style={draw=green!45!black, line width=0.55pt},
      force/.style={-{Latex[length=2.4mm]}, draw=red!75!black, line width=1.2pt},
      disp/.style={-{Latex[length=1.8mm]}, draw=black!45, dashed},
      title/.style={font=\small\bfseries, align=center},
      lab/.style={font=\scriptsize, align=center}
    ]

    \begin{scope}[shift={(0,0)}]
      \node[title] at (0,2.30) {등방 강성\\$k_x=k_y$};
      \draw[axis] (-2.20,0) -- (2.35,0) node[right, lab] {$e_x$};
      \draw[axis] (0,-1.80) -- (0,1.95) node[pos=0.92, left, lab] {$e_y$};
      \foreach \r in {0.55,1.00,1.45} {
        \draw[contour] (0,0) ellipse ({\r} and {\r});
      }
      \node[lab, text=blue!55!black] at (-1.58,1.42) {같은 $U$\\등고선};
      \fill[black] (0,0) circle (1.8pt);
      \node[lab] at (-0.65,-0.42) {목표점\\$\bm{x}_d$, $U_{\min}$};
      \coordinate (piso) at (1.18,0.68);
      \fill[red!75!black] (piso) circle (2.0pt);
      \node[lab] at (1.62,1.18) {현재 오차\\$\bm{e}=\bm{x}-\bm{x}_d$};
      \draw[disp] (0,0) -- node[below, lab] {$\bm{e}$} (piso);
      \draw[force] (piso) -- ++(-0.88,-0.50);
      \node[lab, text=red!75!black] at (0.72,0.28) {$\bm{f}_k=-k\bm{e}$};
      \node[lab, text=black!70] at (0,-2.20) {원형 등고선: 어느 방향으로 밀어도\\같은 강성 기울기를 느낀다.};
    \end{scope}

    \begin{scope}[shift={(6.45,0)}]
      \node[title] at (0,2.30) {방향별 강성\\$k_x<k_y$};
      \draw[axis] (-2.60,0) -- (2.70,0) node[right, lab] {$e_x$};
      \draw[axis] (0,-1.80) -- (0,1.95) node[pos=0.92, left, lab] {$e_y$};
      \foreach \xrad/\yrad in {0.90/0.38,1.55/0.68,2.18/0.96} {
        \draw[softcontour] (0,0) ellipse ({\xrad} and {\yrad});
      }
      \node[lab, text=green!35!black] at (-1.78,1.28) {같은 $U$\\타원 등고선};
      \node[lab, text=green!35!black] at (1.55,-0.42) {$e_x$ 방향\\부드러움};
      \node[lab, text=green!35!black] at (-0.75,1.48) {$e_y$ 방향\\단단함};
      \fill[black] (0,0) circle (1.8pt);
      \node[lab] at (-0.52,-0.42) {목표점\\$\bm{x}_d$};
      \coordinate (pani) at (1.45,0.63);
      \fill[red!75!black] (pani) circle (2.0pt);
      \node[lab] at (2.02,0.92) {현재 오차\\$\bm{e}$};
      \draw[disp] (0,0) -- node[above, lab] {$\bm{e}$} (pani);
      \draw[force] (pani) -- ++(-0.48,-0.92);
      \node[lab, text=red!75!black] at (0.70,-0.88) {$-\bm{K}_d\bm{e}$\\음의 기울기};
      \node[lab, text=black!70] at (0,-2.20) {타원이 긴 방향은 부드럽고,\\짧은 방향은 더 단단하다.};
    \end{scope}
  \end{tikzpicture}%
  }
  \caption{가상 퍼텐셜 에너지와 방향별 강성. 등고선은 같은 힘이 아니라 같은 가상 퍼텐셜 에너지 $U=\frac{1}{2}\bm{e}^{\mathsf T}\bm{K}_d\bm{e}$를 뜻한다. 등방 강성에서는 목표점 주변의 에너지 등고선이 원형이고, 복원력이 변위의 반대 방향과 나란하다. 방향별 강성이 다르면 등고선이 타원형이 되며, 긴 축 방향은 더 부드럽고 짧은 축 방향은 더 단단하다. 로봇 끝단에 작용하는 가상 스프링 복원력은 $\bm{f}_k=-\nabla U=-\bm{K}_d\bm{e}$이고, 끝단을 그 위치에 붙잡기 위해 외부에서 가해야 하는 힘은 반대 방향이다.}
  \label{fig:virtual-potential-stiffness}
\end{figure}


앞에서 같은 $10~\mathrm{N}$이라도 $k=1000~\mathrm{N/m}$에서 $1~\mathrm{cm}$ 밀린 경우와 $k=100~\mathrm{N/m}$에서 $10~\mathrm{cm}$ 밀린 경우가 다르다고 했다.
그 차이는 그림의 에너지 지형으로도 볼 수 있다.
전자는 좁고 가파른 그릇의 벽을 조금 올라간 상태이고, 후자는 완만한 그릇을 훨씬 멀리 이동한 상태이다.
순간 힘은 같을 수 있지만, 저장된 에너지와 주변 기울기, 그리고 사용자가 느끼는 조작감은 다르다.
여기서 말하는 에너지는 실제 금속 스프링 안에 저장된 에너지라기보다, 제어기가 로봇 끝단에 만들고자 하는 가상의 에너지 지형이다.
하지만 가상일과 일률 관점에서는 실제 접촉 거동을 이해하는 데 매우 유용한 해석이다.

감쇠항은 에너지를 줄이는 역할을 한다.
감쇠는 위치만으로 정해지는 보존 퍼텐셜 힘이 아니라, 속도에 의존하는 비보존 소산항이다.
희망 속도가 0이라고 하면 감쇠력은
\begin{equation}
  \bm{f}_d=-\bm{D}_d\dot{\bm{x}}
\end{equation}
이고, 이 감쇠력이 시스템에서 빼앗는 일률은
\begin{equation}
  p_d
  =
  -\bm{f}_d^{\mathsf{T}}\dot{\bm{x}}
  =
  \dot{\bm{x}}^{\mathsf{T}}\bm{D}_d\dot{\bm{x}}.
\end{equation}
즉, 감쇠력이 로봇에 실제로 하는 일률은 $\bm{f}_d^{\mathsf{T}}\dot{\bm{x}}=-\dot{\bm{x}}^{\mathsf{T}}\bm{D}_d\dot{\bm{x}}$이고, 여기서 $p_d$는 양수로 정의한 소산률이다.
$p_d$는 로봇에 새로 들어오는 에너지 속도가 아니라, 양수로 세기로 약속한 빠져나간 에너지의 속도라고 읽으면 된다.
튜토리얼의 단순한 대각 감쇠 모델에서는 보통 $\bm{D}_d=\bm{D}_d^{\mathsf{T}}\succeq0$로 두므로 $p_d\ge0$이고, 감쇠는 에너지를 넣기보다 소산한다.
더 일반적인 행렬에서는 대칭 부분 $(\bm{D}_d+\bm{D}_d^{\mathsf{T}})/2$가 양의 준정부호인지 확인하면 된다.
따라서 강성만 크게 설정하면 목표점으로 강하게 끌려가지만 진동이 커질 수 있고, 감쇠를 함께 설정해야 에너지가 빠져나가며 안정적으로 수렴한다.

\subsection{가상 벽}

임피던스 제어를 설명할 때 좋은 예시는 가상 벽(virtual wall)이다.
지금까지의 부호 규약, 강성, 감쇠, 정상상태 직관을 모두 모아 보는 예제라고 생각하면 된다.
가상 벽은 벽을 넘어간 침투 깊이에만 작용하는 한쪽 방향(one-sided) 스프링-댐퍼이다.
여기서는 $+x$ 방향으로 벽 안쪽에 들어간다고 약속하자.
로봇 끝단의 $x$좌표가 벽 위치 $x_{\mathrm{wall}}$을 넘어가면 침투 깊이를
\begin{equation}
  \delta = \max(0, x - x_{\mathrm{wall}})
  \label{eq:virtual-wall-penetration}
\end{equation}
로 정의할 수 있다.
벽이 고정되어 있고 이미 침투한 상태 $x>x_{\mathrm{wall}}$라면 $\delta=x-x_{\mathrm{wall}}$이므로 $\dot{\delta}=\dot{x}$이다.
즉, $x>x_{\mathrm{wall}}$이고 $\dot{\delta}>0$이면 침투가 더 깊어지는 중이고, $x>x_{\mathrm{wall}}$이지만 $\dot{\delta}<0$이면 이미 들어간 상태에서 빠져나오는 중이다.
따라서 침투한 상태에서의 접근 방향 침투 속도를
\begin{equation}
  \dot{\delta}_+ = \max(0,\dot{\delta})
  \label{eq:virtual-wall-approach-speed}
\end{equation}
로 두자.
아직 침투하지 않은 $x\le x_{\mathrm{wall}}$에서는 벽 힘 전체를 0으로 두므로 이 감쇠항도 작동하지 않는다.
벽이 스프링-댐퍼처럼 작용한다고 하면 가상 벽 모델이 계산해 로봇에 가한다고 보는 힘은
\begin{equation}
  f_{\mathrm{wall}}
  =
  \begin{cases}
    0, & x \le x_{\mathrm{wall}}, \\
    -k_{\mathrm{wall}}\delta
    -
    b_{\mathrm{wall}}\dot{\delta}_+, & x > x_{\mathrm{wall}}.
  \end{cases}
  \label{eq:virtual-wall-force}
\end{equation}
로 모델링할 수 있다.
절차로 읽으면 기준이 훨씬 선명하다.
벽 바깥에서는 아직 접촉하지 않았으므로 $\delta=0$이고 힘도 없다.
끝단이 벽 안쪽으로 파고드는 동안에는 스프링항이 침투 깊이에 비례해 밀어내고, 감쇠항은 파고드는 속도 $\dot{\delta}_+$에만 저항한다.
감쇠는 파고들 때만 켜지고, 빠져나올 때는 벽 안쪽으로 당기지 않는다.
그래서 벽에 닿는 순간의 튕김은 줄이되, 이미 빠져나오는 끝단을 다시 벽 안쪽으로 잡아당기지는 않는다.
일률로 확인해도 같은 결론이 나온다.
침투가 깊어지는 동안 $\dot{\delta}>0$이면 감쇠항 $-b_{\mathrm{wall}}\dot{\delta}_+$와 속도 방향은 반대이므로 벽 감쇠는 에너지를 빼앗는다.
반대로 후퇴 중에는 $\dot{\delta}<0$이고 $\dot{\delta}_+=0$이므로 감쇠항을 꺼 두어 벽이 로봇을 안쪽으로 끌어당기며 에너지를 넣는 상황을 피한다.
이 점이 단순히 $-b_{\mathrm{wall}}\dot{x}$를 항상 더하는 모델과의 차이이다.
다만 $\max$가 들어간 한쪽 방향 모델이므로 접촉 경계와 $\dot{\delta}=0$ 근처에서는 매끄럽지 않을 수 있다.
실제 데모에서는 시간 간격, 포화, 데드밴드, 히스테리시스, 평활화(smoothing), 힘과 속도 제한을 함께 확인해야 한다.
이 모델은 교육용 접근 방향 감쇠이지, 모든 접촉 현상을 설명하는 일반 Kelvin--Voigt 접촉 모델은 아니다.
여기서 $k_{\mathrm{wall}}$이 크면 단단한 벽, 작으면 부드러운 벽이 된다.
$b_{\mathrm{wall}}$은 벽 안쪽으로 더 파고드는 속도, 즉 양의 침투 속도에 대해서만 작용한다.
벽이 움직이거나 벽의 법선 방향이 $x$축과 다르면 $\dot{x}$ 대신 벽에 대한 상대 법선 속도를 써야 한다.
이 글에서는 가장 단순한 고정 평면 벽을 가정해 침투 깊이와 침투 속도를 한 축에서 설명한다.

그림~\ref{fig:virtual-wall-comparison}은 이 관계를 세 가지 관점에서 요약한다.
첫 번째는 벽 위치와 침투 깊이의 기하학적 의미이고, 두 번째는 감쇠가 침투 응답의 튕김을 줄이는 모습이며, 세 번째는 정지 상태에서 힘 크기와 침투 깊이의 관계이다.
여기서 주의할 점은 오른쪽 그래프가 부호 있는 $f_{\mathrm{wall}}$이 아니라 힘의 크기 $|f_{\mathrm{wall}}|$를 그린다는 것이다.
부호 있는 힘은 벽이 로봇을 $-x$ 방향으로 밀기 때문에 음수이고, 처음 볼 때는 힘의 크기를 먼저 읽는 편이 더 직관적이다.
정지 상태, 즉 $\dot{x}=0$에서는
\begin{equation}
  |f_{\mathrm{wall}}|
  =
  k_{\mathrm{wall}}\delta
\end{equation}
이므로 힘-침투 그래프의 기울기가 벽 강성이다.
반면 감쇠항은 이 직선의 기울기를 바꾸는 파라미터가 아니라, 침투 중 속도에 비례해 추가 저항을 만드는 파라미터이다.
따라서 벽에 닿는 순간의 튕김이나 진동을 줄이지만, 완전히 정지한 뒤에는 $\dot{x}=0$이므로 감쇠 힘은 사라진다.

\begin{figure}[!htbp]
  \centering
  \resizebox{\linewidth}{!}{%
  \begin{tikzpicture}[
      axis/.style={-{Latex[length=2.0mm]}, draw=black!55},
      wall/.style={draw=black!70, line width=1.1pt},
      force/.style={-{Latex[length=2.2mm]}, draw=red!75!black, line width=1.2pt},
      motion/.style={-{Latex[length=1.9mm]}, draw=blue!55!black, line width=1.0pt},
      guide/.style={draw=black!40, dashed},
      soft/.style={draw=blue!60!black, line width=1.0pt},
      stiff/.style={draw=red!70!black, line width=1.0pt},
      damped/.style={draw=green!45!black, line width=1.0pt},
      lab/.style={font=\scriptsize, align=center},
      title/.style={font=\small\bfseries, align=center},
      note/.style={font=\scriptsize, align=center, text=black!70}
    ]

    \begin{scope}[shift={(0,0)}]
      \node[title] at (0,2.25) {1D 가상 벽 단면};
      \draw[axis] (-2.05,0) -- (2.35,0) node[right, lab] {$+x$};
      \draw[wall] (0,-1.15) -- (0,1.35);
      \node[lab] at (0,-1.45) {$x_{\mathrm{wall}}$};
      \fill[black!10] (0,-1.15) rectangle (2.0,1.35);
      \node[lab, text=black!65] at (1.20,1.05) {벽 안쪽};
      \node[lab, text=black!65] at (-1.25,1.05) {자유 공간};
      \coordinate (tip) at (1.18,0);
      \fill[blue!55!black] (tip) circle (2.2pt);
      \node[lab] at (1.20,-0.35) {끝단 위치\\$x$};
      \draw[guide] (tip) -- (1.18,0.78);
      \draw[<->, draw=black!55] (0,0.62) -- node[above, lab] {$\delta=x-x_{\mathrm{wall}}$} (1.18,0.62);
      \draw[force] (tip) -- ++(-0.88,0)
        node[midway, below, lab, text=red!75!black] {$f_{\mathrm{wall}}<0$};
      \draw[motion] (0.45,-0.85) -- ++(0.65,0)
        node[midway, below, lab, text=blue!55!black] {침투 중\\$\dot{x}>0$};
      \node[note] at (0,-2.00) {$\delta=\max(0,x-x_{\mathrm{wall}})$\\가상 스프링을 계산하는 침투 거리};
    \end{scope}

    \begin{scope}[shift={(5.4,0)}]
      \node[title] at (0,2.25) {침투 시간 응답};
      \draw[axis] (-0.15,0) -- (3.45,0) node[right, lab] {$t$};
      \draw[axis] (0,-0.15) -- (0,1.55) node[above, lab] {$\delta(t)$};
      \draw[guide] (0,1.0) -- (3.25,1.0);
      \node[lab] at (3.05,1.20) {목표 침투};
      \draw[soft]
        plot[smooth] coordinates {(0,0) (0.45,0.45) (0.85,1.20) (1.25,0.72) (1.70,1.12) (2.15,0.88) (2.70,1.02) (3.15,1.00)};
      \draw[damped]
        plot[smooth] coordinates {(0,0) (0.45,0.55) (0.90,0.90) (1.35,1.02) (1.85,1.01) (2.45,1.00) (3.15,1.00)};
      \node[lab, text=blue!60!black] at (2.65,0.63) {감쇠 작음\\튕김 큼};
      \node[lab, text=green!45!black] at (1.60,1.35) {감쇠 큼\\진동 감소};
      \node[note] at (1.65,-0.58) {감쇠는 기울기를 바꾸기보다\\침투 속도에 브레이크를 건다.};
    \end{scope}

    \begin{scope}[shift={(10.15,0)}]
      \node[title] at (0,2.25) {힘-침투 관계};
      \draw[axis] (-0.15,0) -- (3.25,0) node[right, lab] {$\delta$};
      \draw[axis] (0,-0.15) -- (0,1.75) node[above, lab] {$|f_{\mathrm{wall}}|$};
      \draw[soft] (0,0) -- (3.0,0.85);
      \draw[stiff] (0,0) -- (2.05,1.55);
      \node[lab, text=blue!60!black] at (2.45,0.56) {부드러운 벽\\작은 $k_{\mathrm{wall}}$};
      \node[lab, text=red!70!black] at (2.15,1.35) {단단한 벽\\큰 $k_{\mathrm{wall}}$};
      \node[lab] at (1.30,-0.48) {$\dot{x}=0$이면\\$|f_{\mathrm{wall}}|=k_{\mathrm{wall}}\delta$};
      \draw[guide] (1.35,0) -- (1.35,1.05);
      \draw[<->, draw=black!55] (1.48,0.38) -- node[right, lab] {같은 $\delta$에서\\더 큰 힘} (1.48,1.02);
    \end{scope}
  \end{tikzpicture}%
  }
  \caption{결정론적 가상 벽의 핵심 관계. 이 글의 부호 규약에서는 $+x$ 방향이 벽 안쪽이므로, 침투 깊이는 $\delta=\max(0,x-x_{\mathrm{wall}})$이고 벽이 로봇에 가하는 부호 있는 힘은 $-x$ 방향이다. 초심자에게 힘-침투 관계를 보일 때는 부호 있는 $f_{\mathrm{wall}}$보다 크기 $|f_{\mathrm{wall}}|$를 그리는 편이 직관적이다. 정지 상태에서는 $|f_{\mathrm{wall}}|=k_{\mathrm{wall}}\delta$이므로 기울기가 벽 강성이고, 감쇠항은 기울기를 바꾸는 값이 아니라 접근 방향 침투 속도 $\dot{\delta}_+>0$에 저항해 튕김과 진동을 줄이는 값이다.}
  \label{fig:virtual-wall-comparison}
\end{figure}


MuJoCo의 접촉 해석기에서 제약력을 읽어 올 수도 있지만, 첫 수업에서는 한 단계 단순한 벽이 더 잘 보인다.
여기서 쓰는 결정론적 가상 벽은 접촉 해석기가 알려준 힘을 그대로 쓰는 방식이 아니라, 끝단이 벽을 얼마나 넘었는지 $\delta$와 얼마나 더 파고드는지 $\dot{\delta}_+$를 보고 우리가 정한 식으로 $f_{\mathrm{wall}}$을 계산하는 방식이다.
따라서 이 데모에서 확인할 것은 MuJoCo 접촉 모델의 정확도라기보다, 같은 벽 위치에서 $k_{\mathrm{wall}}$와 $b_{\mathrm{wall}}$을 바꿀 때 침투 깊이, 벽 힘, 끝단 궤적이 어떤 패턴으로 달라지는가이다.

\subsection{실패에서 배우기: 시뮬레이터로 미리 겪는 사고}
\label{sec:failure-gallery}

임피던스 파라미터를 잘못 잡는 실수는 교과서 속 이야기가 아니라, 실제 로봇 앞에서 전공자도 겪는 일이다.
실물에서는 이런 시행착오 한 번이 고가 장비의 파손으로 이어질 수 있지만, 시뮬레이터에서는 같은 사고를 공짜로, 몇 번이고 안전하게 겪어볼 수 있다.
\vmark{fail-gallery-intro}%
이 절의 사례는 저자가 실제 로봇에서 겪은 실수를 단순화한 것이다.
각 사례는 같은 순서로 읽는다.
무슨 일이 벌어지는가, 왜 그런가, 시뮬레이터에서 어떻게 재현하는가, 플롯에서 무엇을 확인하는가이다.

\paragraph{사례: 저장된 탄성 에너지를 계산하지 않고 명령했다.}
목표점을 현재 위치에서 멀리 잡고 강성도 크게 잡았더니, 명령을 넣는 순간 로봇이 목표점 쪽으로 확 튀어나갔다.
이 사고의 원인은 이 장에서 이미 본 가상 퍼텐셜 에너지 식이 전부 설명한다.
목표점에서 $\delta$만큼 떨어진 상태로 강성 $k_d$를 켜는 순간, 가상 스프링에는
\[
  U=\frac{1}{2}k_d\delta^2
\vmark{fail-f2-energy-precheck}%
\]
의 에너지가 이미 장전되고, 손끝에는 최대 $f_{\max}=k_d\delta$의 힘이 걸린다.
이 에너지가 어디로 가는지는 에너지 보존으로 한 줄이면 나온다.
감쇠가 작아 에너지 손실을 잠시 무시하면, 저장된 에너지가 전부 운동 에너지로 바뀌는 순간 $\frac{1}{2}mv^2=\frac{1}{2}k_d\delta^2$이므로 최고 속도는
\[
  v_{\mathrm{peak}}
  \approx
  \delta\sqrt{\frac{k_d}{m}}
\]
이 된다.
즉 튀어나가는 속도는 이동거리와 강성 제곱근의 곱으로, 명령을 넣기 전에 암산할 수 있는 값이다.

숫자로 보면 위험이 선명해진다.
질량 $1~\mathrm{kg}$ 슬라이더를 목표점에서 $0.5~\mathrm{m}$ 떨어뜨려 놓고 강성을 $400~\mathrm{N/m}$로 켜면, 장전 에너지는 $\frac{1}{2}\times400\times0.5^2=50~\mathrm{J}$, 초기 힘은 $200~\mathrm{N}$, 예상 최고 속도는 $0.5\sqrt{400/1}=10~\mathrm{m/s}$이다.
실제로 Lab01에서 이 설정을 실행하면 최고 속도 $9.8~\mathrm{m/s}$가 측정되고(작은 감쇠 $0.5~\mathrm{N\,s/m}$ 때문에 이론값보다 약간 낮다), 4초가 지나도 시스템은 $3.7~\mathrm{m/s}$로 여전히 폭주 중이다.
같은 $0.5~\mathrm{m}$를 강성 $30~\mathrm{N/m}$와 충분한 감쇠 $6~\mathrm{N\,s/m}$로 명령하면 장전 에너지 $3.75~\mathrm{J}$, 초기 힘 $15~\mathrm{N}$, 측정 최고 속도 $1.4~\mathrm{m/s}$로 부드럽게 정착한다.
같은 이동거리인데 설정만으로 장전 에너지가 13배 차이 난다.
\vmark{fail-f2-numeric-contrast}%

두 상황은 저장소의 다음 설정으로 직접 재현할 수 있다.
\begin{itemize}
  \item 위험판: \texttt{configs/lab01\_msd/f2\_launch\_high\_energy.yaml} --- 위치와 속도 플롯에서 명령 직후의 속도 피크와 오래 남는 진동을 확인한다.
  \item 안전판: \texttt{configs/lab01\_msd/f2\_launch\_precheck.yaml} --- 같은 이동거리가 낮은 강성과 충분한 감쇠에서 어떻게 정착하는지 비교한다.
\end{itemize}
\vmark{fail-f2-repro-configs}%
여기서 얻을 습관은 하나다.
큰 목표 오프셋을 명령하기 전에 $\frac{1}{2}k_d\delta^2$과 $k_d\delta$를 먼저 계산해 보고, 그 에너지와 힘이 지금 손끝 주변에서 풀려도 되는 양인지 묻는 것이다.
실물이라면 이 검산 한 줄이 장비 값이다.

\paragraph{사례: 이동 명령을 줬는데 그만큼 가지 않고 멈췄다.}
임피던스 모드에서 목표를 한 방향으로 옮겼는데, 로봇이 중간에 멈춰 서서 목표까지 가지 않았다.
처음에는 고장이나 설정 실수로 보이지만, 이것은 임피던스 제어가 설계대로 동작한 결과일 수 있다.
\vmark{fail-f3-equilibrium-narrative}%
이 절 앞에서 본 평형점 구분이 그대로 답이다.
임피던스 제어의 $x_d$는 반드시 도달해야 하는 위치 명령이 아니라 가상 스프링의 anchor이고, 이동 경로에 접촉이 생기면 로봇은 식 \eqref{eq:steady_position_under_force}의 하중 평형점 $x_{\mathrm{ss}}=x_d+f_{\mathrm{ext}}/k_d$에서 멈춘다.
멈춘 위치와 목표 사이의 간극은 오류가 아니라, 그 간극에 강성을 곱한 만큼의 힘으로 로봇이 환경을 밀고 있다는 표시이다.
또 관절 임피던스를 쓰면서 손끝 직선 이동을 기대했다면, 좌표계 선택부터 다시 볼 일이다.

Lab04의 가상 벽 시나리오가 이 장면을 그대로 재현한다.
목표를 벽 너머로 계속 밀어 넣으면, 손끝은 벽면을 약 $3.3~\mathrm{cm}$ 침투한 지점에서 멈추고 목표만 벽 안쪽으로 앞서간다.
그 순간 벽이 로봇에 가하는 힘은 약 $8.6~\mathrm{N}$인데, 이는 벽 강성 $260~\mathrm{N/m}$에 침투 깊이를 곱한 값과 정확히 일치한다: $260\times0.033\approx8.6~\mathrm{N}$.
즉 로봇은 목표 도달에 실패한 것이 아니라, 명령 스프링의 당김과 벽 힘이 균형을 이룬 위치에서 설계대로 힘을 내고 있는 것이다.
\vmark{fail-f3-wall-balance-numeric}%
재현은 \texttt{configs/lab04\_panda/impedance\_wall.yaml}을 실행해 목표 위치, 손끝 위치, 벽 위치, 계산된 벽 힘 플롯을 함께 읽으면 된다.
\vmark{fail-f3-repro-config}%
여기서 얻을 습관도 하나다.
임피던스 제어에서 로봇이 왜 안 가지라고 느껴지면, 먼저 지금 어떤 외력과 균형 중인가를 묻고, 목표-실측 간극에 강성을 곱해 접촉힘을 추정해 보는 것이다.

\paragraph{사례: 파라미터를 임의로 바꿔 가며 배우려 했다.}
임피던스 제어가 낯설 때 가장 자연스러운 충동은 파라미터를 이것저것 바꿔 보며 감을 잡는 것이다.
시뮬레이터에서는 이 충동이 미덕이지만, 고가의 실물 로봇에서는 한 번의 시도마다 비용과 위험이 붙는 도박이 된다.
\vmark{fail-f1-sweep-narrative}%
게다가 여러 파라미터를 한꺼번에 바꾸면 배움도 사라진다.
강성, 감쇠계수, 질량은 \(\zeta=d_d/(2\sqrt{m_dk_d})\)처럼 서로 묶여 있어서, 두 개를 같이 바꾸면 응답이 달라져도 어느 쪽 때문인지 알 수 없다.
그래서 이 사고의 백신은 앞의 목적별 튜닝 절에서 본 순서 그대로다.
안전 한계를 먼저 정하고, 작은 값에서 시작해, 한 번에 하나만 바꾸고, 매 실행의 플롯을 예측과 비교하는 것이다.
바꾸기 전에 어느 방향으로 변할지 한 줄로 적어 두면, 맞았을 때는 이해가 쌓이고 틀렸을 때는 질문이 생긴다.
\vmark{fail-f1-one-factor-method}%
이 연습을 위한 비교 쌍은 저장소에 이미 준비되어 있다.
Lab01의 \texttt{underdamped}/\texttt{over\_damped}와 \texttt{high\_stiffness}/\texttt{low\_stiffness}는 같은 시스템에서 한 축씩만 바꾼 대비 세트이고, Lab04의 \texttt{wall\_low\_damping}/\texttt{wall\_high\_damping}은 같은 벽 접촉에서 감쇠만 바꾼 쌍이다.
실물 앞에 서기 전에, 이 쌍들로 한 축씩 바꾸는 감각을 공짜로 만들 수 있다.
\vmark{fail-f1-config-pairs}%

이 장을 읽은 뒤에는 다음 네 가지를 구분하면 충분하다.
\begin{itemize}
  \item 이상적인 임피던스 관계는 원하는 힘-운동 관계이고, 실제 명령 힘이나 모터 토크는 그 관계에 가까워지도록 계산한 구현 신호이다. Lab04에서는 이 대신 계산된 가상 벽 힘, 목표 후퇴량, actuator effort proxy를 함께 읽는다.
  \item 정상상태에서는 속도와 가속도가 0에 가까워져 강성과 평형점이 주로 보이고, 과도응답에서는 감쇠와 관성의 영향이 함께 나타난다.
  \item 가상 벽 힘은 좌표와 속도로 직접 계산한 교육용 힘이고, MuJoCo 접촉 해석기의 실제 접촉 제약력과는 구분해서 읽어야 한다.
  \item 플롯에서는 위치 오차, 속도, 계산된 가상 벽 힘, 침투 깊이, 목표 후퇴량, 구동기 부담 지표를 함께 보아야 파라미터 변화의 의미가 보인다.
\end{itemize}

다음 장의 Lab들은 새 이론을 추가하려는 장이 아니라, 지금까지의 힘-운동 관계가 설정값, 로그, 플롯에서 어떻게 보이는지 확인하기 위한 관찰 장치이다.
특히 강성, 감쇠, 목표 오프셋, DLS 오프셋, effort proxy가 서로 다른 층위의 값이라는 점을 실제 실행 결과로 다시 확인한다.
